\documentclass[11pt]{amsart}
\usepackage[T1]{fontenc}
\usepackage{lmodern,amsmath,amssymb,mathtools,microtype,booktabs}
\usepackage[margin=1.15in]{geometry}
\usepackage{xr-hyper}
\usepackage[hidelinks]{hyperref}
\newtheorem{theorem}{Theorem}[section]
\newtheorem{proposition}[theorem]{Proposition}
\newtheorem{lemma}[theorem]{Lemma}
\newtheorem{corollary}[theorem]{Corollary}
\theoremstyle{definition}\newtheorem{definition}[theorem]{Definition}
\theoremstyle{remark}\newtheorem{remark}[theorem]{Remark}
\newcommand{\R}{\mathbb R}
\newcommand{\Z}{\mathbb Z}
\newcommand{\E}{\mathbb E}
\newcommand{\Prob}{\mathbb P}
\newcommand{\one}{\mathbf 1}
\newcommand{\dd}{\,\mathrm d}
\DeclareMathOperator{\dist}{dist}
\DeclareMathOperator{\covol}{covol}
\DeclareMathOperator{\area}{area}
\DeclareMathOperator{\diam}{diam}
\numberwithin{equation}{section}

\externaldocument[H-][nocite]{companion}[A2-v41-companion.pdf]
\title[Scalar collision laws]{Scalar collision laws and recognition of periodic dispersing billiards}
\author{Qian Qi}
\date{October 4, 2026}
\subjclass[2020]{37D50, 52A10, 60J10, 62G05, 68Q32}
\keywords{Dispersing billiards, collision probabilities, inverse geometry,
unknown launch law, support measures, period recognition, finite observations}
\hypersetup{pdftitle={Scalar collision laws and recognition of periodic dispersing billiards},pdfauthor={Qian Qi}}
\begin{document}
\raggedbottom
\begin{abstract}
Two fixed opposite forward collision fields determine a separated
configuration of strictly convex planar obstacles and its unknown
stationary launch probability, up to a common translation. The command
length is positive and fixed; the datum consists of the two spatial
mean fields on the plane. The command length plus the footprint
diameter is smaller than the obstacle separation.
The launch law may be singular or atomic,
and its compact convex support need not have interior. A finite prefix
identity recovers occupation. Three observable support functions then
cancel the unknown footprint and give the difference between an
obstacle support and a contact-chord support. The negative atoms of
the resulting signed curvature measure recover the chord, hence the
obstacles and the complete law. The fields determine every forward
response and the full period group. This exact identification uses
Fourier uniqueness. Finite recovery is obtained separately by support
acquisition and a positive moment reconstruction: on quantitative smooth
classes, the same two commands give polynomial-cost geometric recovery
and an explicit exponential sufficient cost for transportation recovery
of the law. They predict every bounded-length response in local spatial
mean. A known variation bound on the launch density gives finite
\(L^1\) density recovery and uniform raw response prediction on bounded
windows. Finite period decisions retain a positive patch margin;
complete finite configurations require a protected complete aperture.
\end{abstract}
\maketitle
\enlargethispage{2pt}
\section{The collision experiment and the results}
\label{sec:two-field-introduction}
The unknown launch position creates a joint inverse problem: one must
recover both the obstacle geometry and the law that blurs each nominal
command. We prove that two fixed opposite forward commands suffice.
Each is observed as a spatial field, and both use the same unknown
stationary law. Their difference recovers occupation by a finite
prefix formula. Their geometric supports then separate the obstacle
from the footprint by an exact cancellation. This gives a complete
inverse, including the probability law and the full translation
period group. The finite construction uses the same two commands.

\subsection{Observations and the exact inverse}
Let $\mathcal O$ be the union of closed obstacles. For a displacement
$b$, the bit $B_b(y)$ is one when $y\notin\mathcal O$ and the
unreflected segment $[y,y+b]$ meets $\mathcal O$. It is zero for a
solid start and for a free miss. A boundary solid start is zero; a
boundary contact reached from a free start counts as a collision.
For nominal center $x$ the hidden start is $x+Z$, where $Z$ has one
unknown law $\mu$. The observed mean is
\[
                 F_b(x)=\int B_b(x+z)\,\dd\mu(z).
\]
All attempted preparations enter the denominator. No free point,
actual start, obstacle label, contact position, angle or time is
observed. Fresh preparations are independent, conditional on the
chosen commands and prior observations. The finite theorems use exact
nominal controls and count their required numerical precision.

Fix $a=te$, with $|e|=1$ and $t>0$. The exact datum is the ordered
pair of whole-plane spatial fields $(F_a,F_{-a})$, supplied almost
everywhere. It consists of two command settings, not two scalar
measurements. Theorem~\ref{thm:two-field-rigidity} applies to nonempty
locally finite configurations of strictly convex bodies with bounded
diameters and a positive separation $d$. Their boundaries need not
be smooth. The law is any compactly supported probability whose
support $A$ is convex, including a point or segment, and
$t+\diam A<d$. The theorem reconstructs
\[
       \mathcal O-s(A),\qquad A-s(A),\qquad
                   (z\mapsto z-s(A))_\#\mu,
\]
Here $n_\phi=(\cos\phi,\sin\phi)$ and the Steiner point is
$s(K)=\pi^{-1}\int_0^{2\pi}h_K(n_\phi)n_\phi\dd\phi$. This common translation is the
entire ambiguity. Corollary~\ref{cor:two-field-fiber} identifies the
full period group and completes every finite-length forward response.
Neither exact conclusion assumes periodicity.

The central identity uses a matched occupation component $P=C-A$
and the two corresponding collision supports $K_+,K_-$. It is
\[
  2h_P(u)-h_{K_+}(u)-h_{K_-}(u)+t|u\cdot e|=h_C(u)-h_{L_C}(u),
\]
where $L_C$ joins the two transverse contact points of $C$. The
unknown support $h_{-A}$ cancels. The angular curvature measure on
the right consists of a nonatomic positive measure and the two
negative atoms of a segment. Its Jordan decomposition determines
$C$ up to translation. This uses the classical planar support-measure
calculus \cite{Schneider,MartinezMaure}; the observable cancellation
and its use for the joint collision inverse are the present argument.
The subsequent division of a compact occupation transform identifies
the law exactly. The finite law theorem instead uses moment estimates
with explicit conditioning.

These two conclusions have distinct quantitative content. The exact
theorem identifies a characteristic function by continuity from a
dense set of frequencies; it requires no lower bound for the obstacle
transform near its zeros. Finite acquisition first estimates geometric
supports, then fits a positive probability with controlled finitely
many moments. Its transportation error and its exponential sufficient
observation cost are proved directly. Exact response completion and
finite prediction are stated separately in their respective data
models and norms.

\subsection{Quantitative classes and losses}
For finite geometric reconstruction put $s=6+\beta$, where
$0<\beta\le1$. A bounded periodic presentation is
\begin{equation}\label{eq:two-field-presentation}
 \mathcal O=\bigcup_{i=1}^r(C_i+L\Z^2),\qquad 1\le r\le r_0.
\end{equation}
The presentation has exactly $r$ component orbits and may be
nonprimitive. Fix known constants with
\begin{equation}\label{eq:two-field-geometric-priors}
 \begin{gathered}
 \|L\|,\|L^{-1}\|\le L_0,\quad |\det L|\ge V_0>0,\quad
 \diam C_i\le D_0,\quad \dist(C,D)\ge d_0>0,\\
 \kappa_-\le\kappa_C\le\kappa_+,\qquad
                 \|p_C^\circ\|_{C^{6,\beta}}\le M_0.
 \end{gathered}
\end{equation}
Here distinct $C,D$ are physical components, $h_C$ is the support
function, $z_C=s(C)$, and $p_C^\circ=h_C-z_C\cdot n_\phi$.
The positive curvature bounds are uniform. The footprint and its
density have the separate priors
\eqref{eq:two-field-finite-priors}: a smooth support with upper and
lower curvature radii and a lower boundary-mass bound. These
quantitative hypotheses are used only in the finite assertions.

We state the finite period margin explicitly. Put $B=1+L_0$,
$Q_R=[-R,R]^2$, and $\Pi=\operatorname{Per}(\mathcal O)$. Define
\begin{align}
 d_*(C,D)&=\|z_C-z_D\|_\infty+
                         \|p_C^\circ-p_D^\circ\|_\infty,\nonumber\\
 \mathcal D(w)&=\max_{C:z_C\in Q_B}\ \max_{\sigma\in\{-1,1\}}
                            \inf_D d_*(C+\sigma w,D).
                   \label{eq:two-field-patch-defect}
\end{align}
The class $\mathcal B_\eta$ consists of these presentations with a
known $\eta>0$ such that
\begin{equation}\label{eq:two-field-patch-margin}
 \begin{gathered}
 w=z_D-z_C,\quad z_C\in Q_{3B},\ z_D\in Q_{5B},\quad w\notin\Pi
                 \quad\Longrightarrow\quad\mathcal D(w)\ge\eta.
 \end{gathered}
\end{equation}
This is a defect of whole-patch translations. It is invariant under
the common centering, as proved in the companion,
Lemma~\ref{H-lem:registration-period-margin}. The geometric loss is
the maximum $C^2$ error of matched support functions in a protected
patch, together with the errors of a matched primitive period basis
and free area. A basis is matched to a true basis; no continuous
choice of a reduced basis is imposed. Primitive orbit counts and
exact rational relations of the recovered period-generating list are
separate discrete outputs. For a finite configuration the loss is
the maximum $C^2$ error of all matched components in its known
complete protected aperture, and no period margin is assumed.

Theorem~\ref{thm:two-field-finite-geometry} recovers geometry to
$O(\nu)$ from finitely many bits with the fixed commands $\pm te_1$,
where $2t+\Delta<d_0$. Its sufficient power is
\[
            Q_{\rm pair}=\frac{(\gamma+9/2)s}{s-2}.
\]
Theorem~\ref{thm:two-field-finite-law} adds recovery of the canonical
probability in $W_1$ at sufficient cost
\[
 \exp(C\varepsilon^{-1}\log(C/\varepsilon))\log^2(C/\delta).
\]
Corollary~\ref{cor:two-field-finite-prediction} gives the corresponding
local spatial $L^1$ prediction of every bounded-length command.
Under a known total-variation bound on the zero-extended density,
Corollary~\ref{cor:two-field-finite-density} gives density $L^1$ error
and uniform raw-mean error on bounded windows. Its sufficient cost
has $\varepsilon^{-2}$ in place of $\varepsilon^{-1}$. The proofs
give these norms and costs directly; exact Fourier identification is
not used as a finite stability estimate.

\subsection{Resources and relation to the retained results}
We distinguish the number $S$ of distinct nominal sites, the number
$J$ of site-and-sign command occurrences, and the number $N$ of
attempted bits, including repetitions, solid starts and free misses.
Thus $S\le J\le N$. The two displacement vectors are fixed; a
required nominal grid is part of the theorem. Controller description
length, deterministic numerical arithmetic, and physical motion or
positioning are separate resources. A fixed nominal aperture becomes
an absolute physical aperture after a coarse bound on the unknown
footprint origin is included. This bound does not register the origin.

The technical supplement retains the directional-germ inverse, its
one-resolved-obstacle extension, all earlier preparation experiments,
and their complete proofs. The new fixed-pair inverse removes both
the angular continuum of commanded directions and the short-length
limit. Its geometric hypotheses do not compare the footprint diameter
with an obstacle width. For the quantitative smooth class, the
sufficient exponent $Q_{\rm pair}$ is one third of the retained
single-law exponent $(3\gamma+27/2)s/(s-2)$. At $\gamma=0$, $s=7$,
these powers are $6.3$ and $18.9$, respectively. These are sufficient
upper bounds for their stated experiments. The sharp power $2.3$
at $s=7$ in the companion concerns the separate fixed, known
uniform-disk law; it supplies no matched lower bound for the present
unknown-law problem.

The geometric hull step uses known nominal centers with positive
bits and a cap-mass argument of the kind underlying random support
estimation \cite{BrunelCaps}. It does not observe a random point of
the unknown obstacle. The law step uses classical moment comparison
and transportation duality \cite{RigolletWeed}, with a recovered and
therefore perturbed convolution factor. The prediction step uses
bounded variation \cite{AFP,Galerne}. Each of these reductions and
its uniform estimates is proved below. The supplement supplies the
radial interpolation and period-locking proofs through the explicit
interface in Section~\ref{sec:finite-interface}. The three main
sections proceed from exact identification, through finite geometry,
to finite law recovery and prediction. Appendix~\ref{sec:two-field-geometric-details}
gives the complete collar, grazing-cap, grid and assignment estimates;
Appendix~\ref{sec:moment-factor-details} gives the positive-factor
and finite probability reconstruction details.

\subsection{Related inverse problems}
\label{sec:theorem-comparison}
The closest comparisons concern geometric factorization and additive
deconvolution with unknown noise. We describe their data and uniqueness
classes explicitly. This also distinguishes the identification supplied
by the collision fields from the classical inversions used after the
geometric factors have been separated.

For planar convex bodies $C,A$, write
\[
 g_{C,A}(x)=|C\cap(A+x)|
          =(\mathbf1_C*\mathbf1_{-A})(x).
\]
Averkov and Bianchi \cite[Theorem~1.1]{AverkovBianchi} prove that
$g_{C,C}$ determines every planar convex body up to translation and
reflection. The corresponding problem with two unknown bodies has
also been solved in substantial classes. Bianchi
\cite[Theorem~1.1]{BianchiCross} determines a pair of planar polygons
from its cross covariogram, apart from specified parallelogram
exceptions and the unavoidable common translations and reflected
interchanges $(C,A)\mapsto(-A,-C)$. For planar bodies of class
$C^8_+$, meaning $C^8$ boundaries with positive curvature,
\cite[Theorem~6.3]{BianchiLaplace} proves uniqueness up to precisely
these two equivalences. Thus simultaneous recovery of two convex
factors is an established geometric inverse problem.

When the launch law is uniform on a full-dimensional convex body $A$,
our isolated occupation field is $|A|^{-1}g_{C,A}$; the normalization
is part of the comparison. Theorem~\ref{thm:two-field-rigidity} starts
with the ordered collision fields, permits an arbitrary compact law
with convex support, and imposes strict convexity only on the obstacle.
The recovered collision supports provide the additional factorization
identity. Its full fiber is common translation, without a reflected
interchange. The geometric conclusions also extend to separated,
locally finite obstacle configurations. On the statistical side,
Bianchi, Gardner and Kiderlen \cite[Theorems~4.10 and~6.4]{BGK} give
strong Hausdorff consistency for reconstruction from noisy covariogram
grids, using a preliminary difference-body reconstruction and their
convex-body uniqueness and noise assumptions. The finite
experiment below acquires Bernoulli collision bits and controls the
resulting geometric and probability errors in that experiment.

Unknown-probe morphology gives another geometric comparison.
Villarrubia's blind tip reconstruction
\cite[Section~5.1, equations~(14)--(15)]{Villarrubia} successively
intersects dilation constraints, starting from a containing tip bound.
The limiting reconstruction is the largest tip consistent with those
constraints; additional images can be incorporated
\cite[Section~7.3.3]{Villarrubia}. This identifies a maximal consistent
probe, with exact recovery governed by the available contacts.
In the present problem the three matched supports determine both
geometric factors, and the occupation values subsequently determine
the probability carried by the footprint.

For passive additive observations $Y=X+\epsilon$, Gassiat, Le Corff
and Leh\'ericy \cite[Theorem~2.1]{GLL} identify both unknown laws,
up to opposite translations of the summands, when the noise has two
independent coordinate blocks and the signal has an entire Laplace
transform of order less than two satisfying their condition (H2).
Writing the entire characteristic transform as $\Phi$, this condition
requires
\[
 \Phi(z_1,\,\cdot\,)\not\equiv0
 \quad\hbox{and}\quad
 \Phi(\,\cdot\,,z_2)\not\equiv0
 \qquad\hbox{for every complex }z_1,z_2.
\]
Their theorem permits singular laws and zeros of the noise transform.
Capitao-Miniconi, Gassiat and Leh\'ericy
\cite[Corollary~2.5]{CGLSupport} give geometric sufficient conditions
for this hypothesis, including compact strictly convex supports.
Their Theorem~5.1 gives $W_p$ risk for the compact signal law of order
$\log\log n/\log n$ on the stated compact-support classes, which
include a quantitative contrast condition in addition to the exact
identification assumptions.

There is a direct overlap with our occupation reduction. If $U_C$ is
uniform on the recovered component $C$ and independent of the launch
displacement $Z$, the normalized occupation density is the law of
$U_C-Z$. When $Z$ has independent coordinate blocks, the preceding
identification theorem applies, since $C$ is strictly convex. Our
launch law need not have this independence. The two active collision
fields supply the signed support identity that separates the factors
for a general planar launch law. Neither the passive additive samples
nor their statistical rates are assumed available in the finite
collision experiment.

The support-measure step has a precise classical antecedent.
Martinez-Maure \cite[Definition~4.4 and Theorem~4.6]{MartinezMaure}
associates a signed length measure to a difference of support
functions and determines that difference, modulo translations, from
its zero-first-moment measure. Such a difference alone does not
determine its two convex summands. Here the observable identity yields
$h_C-h_{L_C}$. Strict convexity makes $S_C$ nonatomic, whereas the
segment contributes exactly two atoms. The Jordan decomposition
therefore separates the summands under the hypotheses of our exact
theorem, including nonsmooth strictly convex obstacles.

Compact support and Fourier zeros also have a substantial
deconvolution literature. Meister \cite[Section~2 and Theorem~1]{MeisterCompact}
recovers a compactly supported univariate density using the known
error transform only near the origin. Delaigle and Meister
\cite[Theorem~4.1]{DelaigleMeister} obtain statistical bounds for
univariate deconvolution with specified oscillating error transforms
and one-sided signal support. In our exact planar inverse the compact
body indicator has an entire Fourier transform, nonzero at zero.
Its real nonzero set is dense, so division there and continuity
determine the launch transform everywhere. This argument allows
nonisolated planar zero sets. The finite result instead uses mixed
moments with explicit conditioning for the estimated body factor.

Moment control in transportation distance is classical; see
Rigollet and Weed \cite[Section~2.3, Theorem~4 and Corollary~1]{RigolletWeed}
for quantitative univariate bounds. The proof below supplies the
bivariate approximation and convolution estimates needed after the
geometry has been estimated, together with their acquisition from
finite collision bits. The prediction conclusions then use standard
variation bounds \cite{AFP,Galerne}. The resulting identification
chain and finite resource bounds concern the same two prescribed
positive-length commands throughout.

\section{Two fixed forward fields determine the table and launch law}
\label{sec:two-field-rigidity}
Fix a unit vector $e$, a length $t>0$, and write $a=te$.
The only commands in this section are $a$ and $-a$, both made with
the same unknown stationary launch law.  In particular, no angular
or length derivative belongs to the datum.  The two observed means
are the bounded measurable functions
\begin{equation}\label{eq:two-field-data}
 F_+(x)=\int_A B_a(x+z)\dd\mu(z),
 \qquad
 F_-(x)=\int_A B_{-a}(x+z)\dd\mu(z),
 \qquad x\in\R^2.
\end{equation}
These are active spatial fields at one fixed positive command
length.  They use the original first-collision bit, including its
zero value at a solid start, and all attempted preparations remain
in the denominator.

Let $\mathcal O=\bigcup_{C\in\mathcal C}C$ be a nonempty, locally
finite configuration of compact convex bodies with nonempty
interiors, diameters at most $D$, and mutual distances at least
$d>0$.  Each body is strictly convex; no boundary smoothness or
positive curvature bound is imposed in this section.
The unknown stationary law $\mu$ is a compactly supported Borel
probability measure whose support $A$ is nonempty and convex, with
$\diam A\le\Delta$.  The support may be lower dimensional and the
law may have atoms or a singular part.
The only size comparison is
\begin{equation}\label{eq:two-field-gap}
                         t+\Delta<d.
\end{equation}
There is no lower obstacle-width assumption, and the position of
$A$ is unrestricted.  No density or smoothness assumption is made
on the law.  Equality of the two data fields may be understood
almost everywhere with respect to planar Lebesgue measure.
As before, $h_K$ denotes the support function of a
compact set $K$, equivalently of its convex hull, and $s(K)$ denotes
the Steiner point when $K$ is nonempty and convex, including the
lower-dimensional cases.  Supports of observed fields below mean
supports of the nonnegative measures $F_\pm(x)\dd x$.

\begin{lemma}[A finite prefix inverse for occupation]
\label{lem:two-field-prefix}
Put
\begin{equation}\label{eq:two-field-increment}
 r(x)=F_+(x)-F_-(x+a),\qquad
 M=\left\lfloor\frac{D+\Delta}{t}\right\rfloor+1.
\end{equation}
The occupation field is given by the finite formula
\begin{equation}\label{eq:two-field-prefix-inverse}
 v(x):=\int_A\one_{\mathcal O}(x+z)\dd\mu(z)
 =\max_{0\le m\le M}
       \left\{-\sum_{k=0}^{m-1}r(x+ka)\right\}.
\end{equation}
The empty sum is zero.  If arbitrary supplied values of $r$ on
these $M$ sites differ by at most $\epsilon$ from the exact values,
the value returned by the same formula differs from $v(x)$ by at
most $M\epsilon$.  In particular, perturbing each of the two mean
fields by at most $\epsilon$ changes the returned value by at most
$2M\epsilon$.  The formula holds pointwise for the physical mean
fields.  If those fields are supplied only as almost-everywhere
equivalence classes, it determines the corresponding class of $v$.
\end{lemma}
\begin{proof}
For every actual start $y$, convexity and $t<d$ give the
endpoint identity
\begin{equation}\label{eq:two-field-bit-identity}
 B_a(y)-B_{-a}(y+a)
                  =\one_{\mathcal O}(y+a)-\one_{\mathcal O}(y).
\end{equation}
Indeed, if precisely one endpoint is solid, the bit pointing toward
it is one and the other is zero.  If both endpoints are free, the
two bits agree.  If both are solid, they lie in the same convex
body and both bits vanish.  The closed-set convention includes
tangencies and boundary endpoints: a boundary solid start is
zero, while a swept boundary reached from a free start counts as
a collision.  Thus the identity remains true on those sets,
which may have positive probability under a singular law.
Integrating against the same stationary probability measure proves
\begin{equation}\label{eq:two-field-occupation-increment}
                         r(x)=v(x+a)-v(x).
\end{equation}
The nominal shift in $F_-(x+a)$ realizes the endpoint comparison
using only the fixed forward command $-a$ with the original law.

The field $v$ is bounded and measurable.  Positivity of the measure
on every neighborhood of every point in $A$ gives
\begin{equation}\label{eq:two-field-occupation-components}
 \operatorname{supp}(v(x)\dd x)=\bigcup_{C\in\mathcal C}P_C,
 \qquad P_C=C+(-A).
\end{equation}
Indeed, each individual occupation measure is the convolution
of $\one_C(x)\dd x$ with the reflected probability measure,
and the support of a convolution of two nonnegative compactly
supported measures is the sum of their supports.  The sets $P_C$
are also the closures of the components of $\{v>0\}$:
every point of $\operatorname{int}(C+(-A))=\operatorname{int}C+(-A)$
has positive occupation, and outside $C+(-A)$ the contribution
of $C$ is zero.  Boundary occupation values need not vanish.
These compact convex component closures have diameters at most
$D+\Delta$ and mutual gaps at least $d-\Delta>t$.
Every sum in \eqref{eq:two-field-prefix-inverse} telescopes to
$v(x)-v(x+ma)$ and is therefore at most $v(x)$.
If $v(x)=0$, the empty sum attains equality.  Otherwise $x$ is in
one component $P_C$.  Since $Mt>D+\Delta$, the ray
$x,x+a,\ldots,x+Ma$ exits that component by time $M$.
At its first exit it cannot land in another component, because
the component gap exceeds $t$.  Its occupation at that point is
zero, so the corresponding prefix attains equality.
Finally, each prefix sum changes by at most $m\epsilon$, and
taking a maximum is nonexpansive in the maximum norm.  This proves
the stability bounds.
\end{proof}

The inverse is a finite directed version of the prefix construction
in the earlier local experiment.  It is useful here because the
two commanded directions are retained as separate channels.
It has a fixed finite spatial range $Mt$ and does not require a
limiting walk, an angularly averaged forcing, or prior knowledge of
the unknown positive components.

The bounded measurability of the fields is sufficient here: each
prefix uses only finitely many fixed translations, so changing
the supplied fields on null sets changes the recovered occupation
only on a null set.

\subsection{Measure supports and singular launch laws}
\label{subsec:measure-support}
For a nonnegative locally integrable function $f$, the convention
used here is
\[
 \operatorname{supp}(f(x)\dd x)
 =\left\{x:\int_{B(x,r)}f(y)\dd y>0
                         \text{ for every }r>0\right\}.
\]
This set depends only on the almost-everywhere class of $f$.
It need not be the pointwise positivity set.  We write
$\check\mu=(-\operatorname{id})_\#\mu$.

\begin{lemma}[Positive convolution supports]
\label{lem:positive-convolution-support}
Let $\alpha,\beta$ be nonzero, nonnegative, compactly supported
finite Borel measures on $\R^2$.  Then
\begin{equation}\label{eq:positive-convolution-support}
 \operatorname{supp}(\alpha*\beta)
            =\operatorname{supp}\alpha+\operatorname{supp}\beta.
\end{equation}
If bounded measurable sets $E,E'$ differ by a planar null set,
then $\one_E*\beta=\one_{E'}*\beta$ almost everywhere,
including when $\beta$ is singular.
\end{lemma}
\begin{proof}
Put $U=\operatorname{supp}\alpha$ and
$V=\operatorname{supp}\beta$.  The convolution is carried by the
compact set $U+V$.  Conversely, if $x=p+q$ with $p\in U$ and
$q\in V$, then, for every $r>0$,
\[
 (\alpha*\beta)(B(x,r))
 \ge \alpha(B(p,r/3))\beta(B(q,r/3))>0.
\]
This proves \eqref{eq:positive-convolution-support}.
Tonelli's theorem identifies the convolution
$(\one_E(x)\dd x)*\beta$ with the measure having density
$\one_E*\beta$.  It also gives
\[
 \int_{\R^2}|\one_E*\beta-\one_{E'}*\beta|\dd x
 \le \beta(\R^2)\,\area(E\mathbin\triangle E')=0.
\]
Neither argument requires a density for either measure.
\end{proof}

\begin{proposition}[Observable components under an arbitrary compact law]
\label{prop:singular-support-components}
Under \eqref{eq:two-field-gap}, let
$v_C(x)=\int\one_C(x+z)\dd\mu(z)$, and let
$f_{\pm,C}(x)=\int B^{C}_{\pm a}(x+z)\dd\mu(z)$,
where $B^{C}$ is the collision bit for the single obstacle $C$.
Use coordinates $y=ue+we^\perp$ and write the sections of $C$ as
$[\ell_C(w),u_C(w)]$, with $w\in I_C=[w_-,w_+]$.
Define the closed arcs and strips by
\begin{align*}
 \Gamma_{+,C}&=\{\ell_C(w)e+we^\perp:w\in I_C\},&
 E_{+,C}&=\Gamma_{+,C}+[-t,0]e,\\
 \Gamma_{-,C}&=\{u_C(w)e+we^\perp:w\in I_C\},&
 E_{-,C}&=\Gamma_{-,C}+[0,t]e.
\end{align*}
Then
\begin{equation}\label{eq:singular-individual-supports}
 \operatorname{supp}(v_C(x)\dd x)=P_C:=C+(-A),\qquad
 \operatorname{supp}(f_{\pm,C}(x)\dd x)
                         =K_{\pm,C}:=E_{\pm,C}+(-A).
\end{equation}
The connected components of the occupation support are precisely
the $P_C$; those of the two collision supports are, respectively,
the $K_{+,C}$ and the $K_{-,C}$.  Moreover,
\begin{align}
 \dist(P_C,P_{C'})&\ge d-\Delta,\label{eq:singular-occupation-gap}\\
 \dist(K_{\pm,C},K_{\pm,C'})&\ge d-t-\Delta,
       \qquad C\ne C',\label{eq:singular-collision-gap}\\
 K_{\pm,C}\cap P_C&\ne\varnothing,\qquad
 \dist(K_{\pm,C},P_{C'})\ge d-t-\Delta
       \quad(C\ne C').\label{eq:singular-support-matching}
\end{align}
Thus intersection with the occupation support determines the
correspondence between the two collision component families.
These statements apply to point, segment, atomic and singular
continuous launch laws with the stated convex support.
\end{proposition}
\begin{proof}
The section functions $\ell_C$ and $u_C$ are convex and concave,
respectively, on $I_C$.  They are continuous in its interior.
Strict convexity makes each extreme section a singleton.  Any
convergent subsequence of section endpoints as $w\to w_\pm$
lies in that singleton, by compactness of $C$.  Both section
functions are consequently continuous on $I_C$.

For the positive command, the single-obstacle collision-start set
is exactly
\[
 D_{+,C}=\{(\ell_C(w)-s)e+we^\perp:
                         w\in I_C,\ 0<s\le t\}.
\]
The endpoint $s=0$ is a solid start and is excluded; $s=t$ is
a collision at the terminal endpoint and is included.  At
$w=w_\pm$ the same formula includes a tangential encounter.
The open set
\[
 G_{+,C}=\{(\ell_C(w)-s)e+we^\perp:
                    w_-<w<w_+,\ 0<s<t\}
\]
has closure $E_{+,C}$.  Indeed, the parametrization is a
homeomorphism from $I_C\times[0,t]$ to $E_{+,C}$, and its
restriction to the open rectangle has open image.  The sets
$D_{+,C}$ and $G_{+,C}$ differ only on two graphs and two
transverse endpoint segments, all of which are planar null sets
by one-dimensional slicing.  The same description, with
$u_C(w)+s$ in place of $\ell_C(w)-s$, applies to the negative
command.  In particular,
\[
 \operatorname{supp}(\one_{G_{\pm,C}}(x)\dd x)=E_{\pm,C}.
\]
Every neighborhood of a point of $E_{\pm,C}$ meets its open
interior in positive area, including at a tangent endpoint.

Since $C$ is the closure of its interior,
$\operatorname{supp}(\one_C(x)\dd x)=C$.  Apply
Lemma~\ref{lem:positive-convolution-support} to these three
Lebesgue measures and to $\check\mu$, whose support is $-A$.
The null-set assertion in that lemma passes from $G_{\pm,C}$
to the physical collision bits after convolution and proves
\eqref{eq:singular-individual-supports}.  This passage is an
almost-everywhere assertion; it does not discard the prescribed
pointwise boundary convention for the physical means.

For completeness, the occupation is positive at every point of
$\operatorname{int}P_C$.  The convex-set identity needed here is
\begin{equation}\label{eq:singular-interior-sum}
             \operatorname{int}(C+(-A))
                         =\operatorname{int}C+(-A).
\end{equation}
To prove the nontrivial inclusion, fix $c_0\in\operatorname{int}C$
and $a_0\in A$.  If $x\in\operatorname{int}(C-A)$, then for
small $\lambda>0$ the point
$(x-\lambda(c_0-a_0))/(1-\lambda)$ belongs to $C-A$.
Write it as $c-a$.  The identity
\[
 x=((1-\lambda)c+\lambda c_0)
                  -((1-\lambda)a+\lambda a_0)
\]
places $x$ in $\operatorname{int}C-A$; the reverse inclusion
follows because the latter set is a union of open translates.
If $x=c-a$ with $c\in\operatorname{int}C$ and $a\in A$, a
sufficiently small neighborhood of $a$ is mapped by $z\mapsto x+z$
into $C$.  That neighborhood has positive $\mu$-mass, so
$v_C(x)>0$.  Outside $P_C$ the occupation is zero.  Its positive
set therefore contains $\operatorname{int}P_C$ and is contained
in $P_C$; it is path connected, since a segment joining any of
its points to an interior point stays in the interior except
possibly at its first endpoint.  Its closure is $P_C$.

Each $E_{\pm,C}$ is path connected by its rectangle
parametrization.  Convexity of $A$ makes $K_{\pm,C}$ path
connected as well.  The estimates
\eqref{eq:singular-occupation-gap}--\eqref{eq:singular-support-matching}
follow by subtracting two representations of their points:
obstacle points have distance at least $d$, the difference of
two same-sign longitudinal shifts has length at most $t$, and
the difference of two footprint shifts has length at most
$\Delta$.  The nonempty intersection follows from
\[
 \Gamma_{\pm,C}+(-A)\subset K_{\pm,C}\cap P_C.
\]

All expanded families remain locally finite: a member meeting a
fixed compact set comes from an obstacle meeting its enlargement
by the fixed bounded set $A+[-t,t]e$.  Since $t<d$, a segment
cannot meet two different obstacles or start in one obstacle and
reach another.  Thus $v=\sum_Cv_C$ and
$F_\pm=\sum_C f_{\pm,C}$, with the physical boundary convention.
Their measure supports are the locally finite unions in
\eqref{eq:singular-individual-supports}.  The positive gaps and
connectedness just proved identify their connected components
and establish the matching assertion.
\end{proof}

For example, for any nonempty compact convex $A$, choose a sequence
$(a_j)_{j\ge1}$ dense in $A$ and put
$\mu=\sum_{j\ge1}2^{-j}\delta_{a_j}$.
Its topological support is $A$, and the proposition applies to this
purely atomic probability. This includes convex supports of dimension
zero, one or two and illustrates why the measure-support convention
is used throughout.

\begin{corollary}[Almost-everywhere data and the period group]
\label{cor:singular-support-ae-periods}
The almost-everywhere classes of $F_+,F_-$ determine the same
occupation and collision support components, with the same
matching, as their physical representatives.  If
\[
 \operatorname{Per}_{\mathrm{ae}}(F_+,F_-)
 =\{w:F_\pm(\cdot+w)=F_\pm\text{ almost everywhere}\},
\]
then
\begin{equation}\label{eq:singular-support-periods}
 \operatorname{Per}_{\mathrm{ae}}(F_+,F_-)
 =\operatorname{Per}_{\mathrm{point}}(F_+,F_-)
 =\operatorname{Per}(\mathcal O).
\end{equation}
\end{corollary}
\begin{proof}
The finite prefix formula in Lemma~\ref{lem:two-field-prefix}
uses only finitely many fixed translations.  A change of the
fields on null sets changes its output only on a finite union
of translated null sets.  It therefore recovers the
almost-everywhere occupation class and hence its measure support.
The two original measure supports and their incidence relation
are likewise unchanged.  Proposition~\ref{prop:singular-support-components}
identifies the components without any choice of representatives.

If $w$ is an almost-everywhere common field period, the same
prefix formula gives $v(\cdot+w)=v$ almost everywhere.
Translation by $w$ permutes the components of its measure
support, so $P_C+w=P_{C'}$ for some $C'$.  Taking support
functions and canceling the common $h_{-A}$ gives
$h_C(u)+w\cdot u=h_{C'}(u)$ for every unit $u$.
Hence $C+w=C'$ and $\mathcal O+w=\mathcal O$.
Conversely, a period of $\mathcal O$ preserves every physical
collision bit pointwise, and integration against the same law
preserves every physical mean pointwise.  This proves
\eqref{eq:singular-support-periods}.
\end{proof}


\begin{lemma}[The two collision supports]
\label{lem:two-field-supports}
For each body $C$, let $\Gamma_{+,C}$ be its incoming boundary arc
for direction $e$, including its two tangent endpoints, and let
$\Gamma_{-,C}$ be the opposite incoming arc for direction $-e$.
Define
\begin{equation}\label{eq:two-field-strips}
 E_{+,C}=\Gamma_{+,C}+[-t,0]e,\qquad
 E_{-,C}=\Gamma_{-,C}+[0,t]e,
 \qquad K_{\pm,C}=E_{\pm,C}+(-A).
\end{equation}
The connected components of $\operatorname{supp}F_\pm$ are exactly
the sets $K_{\pm,C}$.  Within either sign, different components
have distance at least $d-\Delta-t>0$.
The component $K_{\pm,C}$ meets $P_C$ and is disjoint from all
$P_{C'}$ with $C'\ne C$.  Thus the known occupation components
match the two families of collision supports without an unknown
label correspondence.  Their support functions satisfy
\begin{align}
 h_{K_{+,C}}(u)
   &=h_{\Gamma_{+,C}}(u)+h_{-A}(u)+t(-u\cdot e)_+,
       \label{eq:two-field-plus-support}\\
 h_{K_{-,C}}(u)
   &=h_{\Gamma_{-,C}}(u)+h_{-A}(u)+t(u\cdot e)_+.
       \label{eq:two-field-minus-support}
\end{align}
\end{lemma}
\begin{proof}
Use coordinates $y=ue+w e^\perp$.  The section of $C$ at transverse
coordinate $w$ is an interval $[\ell_C(w),u_C(w)]$ for
$w$ in its projection interval $[w_-,w_+]$.
The incoming graph $\ell_C$ is convex and continuous, and the
outgoing graph $u_C$ is concave and continuous.  The free starts
whose segment first hits $C$ in direction $e$ form, up to null
sets, the strip
\[
 \{we^\perp+(\ell_C(w)-s)e:
                         w_-<w<w_+,\ 0<s<t\}.
\]
Its closure is $E_{+,C}$; the reverse direction gives $E_{-,C}$.
No such strip meets another obstacle because $t<d$.
The planar density of $F_\pm(x)\dd x$ is consequently the locally
finite sum of convolutions of the strip indicators with the
reflected measure.  The strips and collision-start sets differ
only on planar null sets; convolution with any finite measure
preserves that equality almost everywhere, by Fubini's theorem.
Each strip is connected, has nonempty interior, and is the
closure of that interior.  The support of its convolution with
the nonnegative measure is exactly its Minkowski sum with $-A$:
there is no cancellation, and both factors give positive mass
to every neighborhood of every point in their supports.
Thus each $K_{\pm,C}$ is a connected compact support.

For two points from different strips with the same sign, the
two longitudinal shifts differ by at most $t$, while two
footprint shifts differ by at most $\Delta$.  Their distance is
therefore at least $d-t-\Delta$.  The same estimate separates
$K_{\pm,C}$ from $P_{C'}$ when $C'\ne C$.
Since $\Gamma_{\pm,C}\subset C$,
$\Gamma_{\pm,C}+(-A)\subset K_{\pm,C}\cap P_C$.
Local finiteness and the positive gap now prove the support
component and matching assertions.  Finally, support functions
add under Minkowski addition, giving the two displayed identities.
\end{proof}

\subsection{Support arcs and angular curvature without boundary regularity}
\label{subsec:nonsmooth-support-curvature}
The incoming arcs can be treated through sections, without a
parametrization by the boundary normal.  This also distinguishes
strict convexity from smoothness in the subsequent measure
calculation.

\begin{lemma}[A section form of the arc identity]
\label{lem:nonsmooth-arc-support}
Let $C$ be a strictly convex planar body with nonempty interior.
Let $p_-,p_+$ be its unique contacts with the support lines of
normals $-e^\perp,e^\perp$, and put $L=[p_-,p_+]$.
For its closed incoming arcs $\Gamma_+,\Gamma_-$ one has
\begin{align}
 h_{\Gamma_+}(u)&=
 \begin{cases}h_C(u),&u\cdot e\le0,\\
                h_L(u),&u\cdot e\ge0,
 \end{cases}
 \label{eq:nonsmooth-plus-hemispheres}\\
 h_{\Gamma_-}(u)&=
 \begin{cases}h_L(u),&u\cdot e\le0,\\
                h_C(u),&u\cdot e\ge0.
 \end{cases}
 \label{eq:nonsmooth-minus-hemispheres}
\end{align}
The two alternatives agree when $u\cdot e=0$.  In particular,
\begin{equation}\label{eq:nonsmooth-arc-identity}
                  h_{\Gamma_+}+h_{\Gamma_-}=h_C+h_L.
\end{equation}
\end{lemma}
\begin{proof}
Write $u=\alpha e+\beta e^\perp$ and use the interval sections
$[\ell(w),r(w)]$ for $w\in[w_-,w_+]$.  Convexity of $C$ makes
$\ell$ convex and $r$ concave.  They are continuous on the
closed projection interval: interior continuity follows from
convexity, and endpoint continuity follows from compactness and
the uniqueness of the extreme sections.  The endpoints of both
graphs are $p_-,p_+$.

If $\alpha\le0$, maximizing $\alpha s+\beta w$ over each section
chooses $s=\ell(w)$.  Consequently
$h_{\Gamma_+}(u)=h_C(u)$.  If $\alpha\ge0$, the function
$w\mapsto\alpha\ell(w)+\beta w$ is convex.  Its value at an
interior point is at most the larger endpoint value, by the
defining convexity inequality.  Its maximum is therefore
$\max\{u\cdot p_-,u\cdot p_+\}=h_L(u)$.
For the other arc, $\alpha\ge0$ selects the right endpoint of
each section and gives $h_C$; when $\alpha\le0$ the function
$\alpha r(w)+\beta w$ is convex and gives $h_L$.
At $\alpha=0$ both maxima depend only on the projection interval.
This proves all the formulas, including directions supporting
a corner and the two transverse directions.
\end{proof}

We identify the angular circle with $\R/(2\pi\Z)$, put
$n_\phi=(\cos\phi,\sin\phi)$ and
$\tau_\phi=n_\phi^\perp$, and use the planar support convention
\begin{equation}\label{eq:nonsmooth-area-convention}
 \int\psi\dd\mathsf S_K
     =\int_0^{2\pi}h_K(n_\phi)
                      (\psi''(\phi)+\psi(\phi))\dd\phi,
 \qquad \psi\in C^\infty(\R/(2\pi\Z)).
\end{equation}
Thus angular differentiation is distributional.  The following
proof supplies both positivity and the precise atom convention
for this surface-area measure.

\begin{lemma}[Surface-area atoms and exposed segments]
\label{lem:nonsmooth-surface-atoms}
For a compact convex planar body $K$ with nonempty interior,
\eqref{eq:nonsmooth-area-convention} defines a finite nonnegative
measure.  Let
\[
 F(K,n_\phi)=\{x\in K:x\cdot n_\phi=h_K(n_\phi)\}
\]
be its exposed face.  Writing $h(\phi)=h_K(n_\phi)$, its
one-sided derivatives satisfy
\begin{equation}\label{eq:nonsmooth-face-derivatives}
 h'_-(\phi)=\min_{x\in F(K,n_\phi)}x\cdot\tau_\phi,
 \qquad
 h'_+(\phi)=\max_{x\in F(K,n_\phi)}x\cdot\tau_\phi.
\end{equation}
Consequently
\begin{equation}\label{eq:nonsmooth-atom-face-length}
 \mathsf S_K(\{n_\phi\})
       =h'_+(\phi)-h'_-(\phi)
       =\operatorname{length}F(K,n_\phi).
\end{equation}
In particular, $\mathsf S_K$ is nonatomic if and only if $K$ is
strictly convex.  Boundary corners and a singular continuous
part of $\mathsf S_K$ are permitted in this equivalence.
\end{lemma}
\begin{proof}
Choose $R$ with $K\subset B(0,R)$.  The support function is
Lipschitz in $\phi$.  Each function
$x\cdot n_\phi+R\phi^2/2$, $x\in K$, is convex on the real
line, since its second derivative is $R-x\cdot n_\phi\ge0$.
Their supremum $h(\phi)+R\phi^2/2$ is convex.  Hence $h$ has
one-sided derivatives of locally bounded variation and its
distributional second derivative is a locally finite signed
measure.  Periodicity gives a finite signed measure on the circle.

If $p$ is a support point at $\phi$, then, for $|b|<\pi/2$,
\[
 h(\phi+b)+h(\phi-b)
 \ge p\cdot(n_{\phi+b}+n_{\phi-b})
 =2\cos b\,h(\phi).
\]
Multiply by a nonnegative smooth periodic $\psi$, integrate,
translate the two shifted integrals, and divide by $b^2$.
Letting $b\to0$ yields
$\int h(\psi''+\psi)\ge0$.  Thus the signed measure
$h''+h\dd\phi$ is nonnegative and is the measure in
\eqref{eq:nonsmooth-area-convention}.

For $b>0$ and $p\in F(K,n_\phi)$, the support inequality gives
\[
 \frac{h(\phi+b)-h(\phi)}{b}
       \ge p\cdot\frac{n_{\phi+b}-n_\phi}{b}.
\]
Conversely, choose $p_b\in F(K,n_{\phi+b})$.  The same
quotient is at most
$p_b\cdot(n_{\phi+b}-n_\phi)/b$.
Every subsequential limit of $p_b$ belongs to $F(K,n_\phi)$,
by compactness and continuity of $h$.  Taking lower and upper
limits proves the right-derivative formula in
\eqref{eq:nonsmooth-face-derivatives}; negative $b$ gives the
left-derivative formula.  The atom of the Stieltjes measure
$h''$ at $\phi$ is the jump $h'_+(\phi)-h'_-(\phi)$, whereas
$h\dd\phi$ has no atoms.  The face is a segment on the line
$x\cdot n_\phi=h(\phi)$, so its length is exactly the difference
between the two extremal $\tau_\phi$ coordinates.  This proves
\eqref{eq:nonsmooth-atom-face-length}.

The geometric interpretation of the measure follows in the same
coordinates.  The right-continuous support contact
$x_+(\phi)=h(\phi)n_\phi+h'_+(\phi)\tau_\phi$ selects the
positive tangent endpoint of each exposed face.  As the normal
turns once, these contacts traverse the convex boundary in order;
each jump is completed by the corresponding exposed segment.
The distributional product rule gives
\[
                      D x_+=\tau_\phi\dd\mathsf S_K.
\]
Since $|\tau_\phi|=1$ and $\mathsf S_K\ge0$, the variation
of this boundary parametrization is $\mathsf S_K$.  It assigns
boundary arclength to its outer normal, with a whole exposed
segment assigned to its fixed normal.  This is the usual planar
surface-area measure in the convention
\eqref{eq:nonsmooth-area-convention}.

A strictly convex body has only singleton exposed faces and
hence no atoms.  Conversely, any nontrivial boundary segment
lies in an exposed face: a supporting line at an interior point
of the segment contains both endpoints.  Such a segment would
give a positive atom.  At a corner of a strictly convex body,
the normal may vary over an interval while the support point
remains a fixed point $p$.  On the interior of that interval
$h(\phi)=p\cdot n_\phi$, so $h''+h=0$ there; its endpoints
also have zero atomic mass by
\eqref{eq:nonsmooth-atom-face-length}.  None of these arguments
excludes a singular continuous part of the measure.
\end{proof}

These are the planar support-measure conventions of
\cite[Chapter~4]{Schneider}.  Signed support differences and the
separate atomic, absolutely continuous and singular continuous
parts of their length measures are treated in
\cite[Section~4]{MartinezMaure}.  Formula
\eqref{eq:nonsmooth-area-convention} retains all three parts;
an almost-everywhere second derivative alone would not do so.

\begin{lemma}[Jordan recovery of the contact chord]
\label{lem:nonsmooth-jordan-chord}
For the body and chord in Lemma~\ref{lem:nonsmooth-arc-support},
put $H=h_C-h_L$, $m=(p_-+p_+)/2$, and
$\ell=|p_+-p_-|$.  If $n$ is either unit normal of the chord,
then the Jordan decomposition of $\sigma_H=(\partial_\phi^2+1)H$
is
\begin{equation}\label{eq:nonsmooth-jordan-parts}
 \sigma_H^+=\mathsf S_C,\qquad
 \sigma_H^-=\ell(\delta_n+\delta_{-n}).
\end{equation}
It determines the centered chord
$L^0=[-\ell n^\perp/2,\ell n^\perp/2]$, and
\begin{equation}\label{eq:nonsmooth-jordan-recovery}
                     H+h_{L^0}=h_{C-m}.
\end{equation}
Exactly one of the two chord normals has positive scalar product
with $e$.
\end{lemma}
\begin{proof}
The projection interval has positive length, so $\ell>0$.
Writing $q=(p_+-p_-)/\ell$, the chord support is
\[
 h_L(n_\phi)=m\cdot n_\phi+\frac\ell2|q\cdot n_\phi|.
\]
The translation term is annihilated by $\partial_\phi^2+1$.
The second term solves $h''+h=0$ away from the two normals
$n,-n$, and its first derivative jumps upward by $\ell$ at
each.  Hence
\[
 \sigma_H=\mathsf S_C-\ell(\delta_n+\delta_{-n}).
\]
Lemma~\ref{lem:nonsmooth-surface-atoms} shows that $\mathsf S_C$
assigns zero mass to these two points.  The two nonnegative
measures on the right are therefore mutually singular, so they
are exactly the positive and negative Jordan parts.  The
locations and equal masses of the negative atoms determine
$L^0$, independent of the choice between $n$ and $-n$.
Since $h_L=h_{L^0}+m\cdot u$, adding $h_{L^0}$ to $H$ gives
\eqref{eq:nonsmooth-jordan-recovery}.  Finally,
$(p_+-p_-)\cdot e^\perp=w_+-w_->0$; a chord normal therefore
has nonzero scalar product with $e$, and its sign selects exactly
one of the two normals.  The chord need not be perpendicular to
the command direction.
\end{proof}


\begin{lemma}[Calibration by a pair of supports]
\label{lem:two-field-calibration}
Let $P=P_C$ and let $K_+,K_-$ be its matched collision components.
The observable function
\begin{equation}\label{eq:two-field-calibration-function}
 H_C(u)=2h_P(u)-h_{K_+}(u)-h_{K_-}(u)+t|u\cdot e|
\end{equation}
determines $C$ up to translation.  More precisely, let $p_+,p_-$
be the two points of $C$ with outward normals $e^\perp,-e^\perp$,
and let $L_C=[p_-,p_+]$.  Then
\begin{equation}\label{eq:two-field-chord-cancellation}
                         H_C=h_C-h_{L_C}.
\end{equation}
Writing $u=n_\phi$, $\ell=|p_+-p_-|>0$, and choosing the unit
normal $n$ of this chord with $n\cdot e>0$, one has, as signed
measures on the angular circle,
\begin{equation}\label{eq:two-field-curvature-measure}
 (\partial_\phi^2+1)H_C
    =\mathsf S_C-\ell(\delta_n+\delta_{-n}),
\end{equation}
Here $\mathsf S_C=h_C''+h_C$ is the planar surface area measure,
which is nonnegative and has no atoms because $C$ is strictly
convex.  The negative part therefore determines the centered segment
$L_C^0=[-\ell n^\perp/2,\ell n^\perp/2]$.  The function
$H_C+h_{L_C^0}$ is the support function of $C-s(L_C)$.
\end{lemma}
\begin{proof}
On the hemisphere $u\cdot e\le0$, the support point of $C$ in
direction $u$ lies on $\Gamma_{+,C}$, so
$h_{\Gamma_{+,C}}(u)=h_C(u)$.  On the opposite hemisphere,
the function
\[
 w\longmapsto (u\cdot e)\ell_C(w)+(u\cdot e^\perp)w
\]
is convex.  Its maximum is attained at a projection endpoint,
so $h_{\Gamma_{+,C}}(u)=h_{L_C}(u)$ there.
For the reverse arc the two hemispheres are interchanged.
Thus, including their endpoints,
\begin{equation}\label{eq:two-field-arc-sum}
              h_{\Gamma_{+,C}}+h_{\Gamma_{-,C}}
                                      =h_C+h_{L_C}.
\end{equation}
Since $h_P=h_C+h_{-A}$, substitution of
\eqref{eq:two-field-plus-support} and
\eqref{eq:two-field-minus-support} into
\eqref{eq:two-field-calibration-function} proves
\eqref{eq:two-field-chord-cancellation}.  In particular, the unknown
footprint cancels exactly before differentiation is used.

For a planar convex body, $h_C''+h_C$ is the finite nonnegative
surface area measure.  Its atom at a normal direction has mass
equal to the length of the exposed boundary segment with that
normal.  In the strictly convex case it therefore has no atoms.
Equivalently, strict convexity gives a unique support point in
every direction; compactness makes that point continuous in the
direction, and differentiation of the support maximum gives
$h_C\in C^1$, so its first derivative has no jumps.  These are
the usual planar support-measure identities; see \cite{Schneider}. Signed differences of planar surface area
measures and their uniqueness up to translation are treated in
\cite[Section~4]{MartinezMaure}. Here the collision identity
\eqref{eq:two-field-chord-cancellation} supplies the particular
signed support whose atomic part can be separated.
A segment of length
$\ell$ has angular curvature measure
$\ell(\delta_n+\delta_{-n})$.  To check the coefficient directly,
after subtracting the midpoint translation its support function is
$\frac\ell2|n^\perp\cdot n_\phi|$: its first derivative has a
jump of $\ell$ at each of the two normal directions and it solves
$h''+h=0$ between them.
This proves \eqref{eq:two-field-curvature-measure}.
The nonnegative nonatomic measure and the two negative atoms are
mutually singular, so they are precisely the positive and negative
parts of this signed measure.  The two endpoints of the segment
have distinct transverse coordinates, hence $\ell>0$ and exactly
one of its normals has positive dot product with $e$.
Its length and unoriented direction are thus recovered from the
negative part.  Subtracting the segment midpoint from
\eqref{eq:two-field-chord-cancellation} gives the last assertion.
\end{proof}

The differentiation in \eqref{eq:two-field-curvature-measure} is
with respect to the support direction of a reconstructed spatial
set.  It introduces no new command direction and no derivative of
the observed response.  The same signed measure also gives the
entire surface area measure as its positive part.  When $C$ has
$C^2$ boundary and positive curvature, this part is
$r_C(\phi)\dd\phi$.

\begin{theorem}[Joint rigidity from two fixed forward fields]
\label{thm:two-field-rigidity}
Under \eqref{eq:two-field-gap}, the two fields
\eqref{eq:two-field-data} determine the complete canonical triple
\begin{equation}\label{eq:two-field-canonical-triple}
  \mathcal O_0=\mathcal O-s(A),\qquad
  A_0=A-s(A),\qquad \mu_0=(z\mapsto z-s(A))_\#\mu,
\end{equation}
with probability measures identified exactly.  In particular,
if $\mu$ has an $L^1$ density $j$, its centered density
$j_0(z)=j(z+s(A))$ is identified almost everywhere.
No angular continuum, short-flight limit, second stationary law,
or separately prescribed reciprocal preparation is required.
No periodicity assumption or lower bound on obstacle widths is
required.
\end{theorem}
\begin{proof}
Recover $v$ by Lemma~\ref{lem:two-field-prefix}, then take the
connected components $P$ of $\operatorname{supp}(v(x)\dd x)$.
Match the two collision
support components by Lemma~\ref{lem:two-field-supports}.
For each $P=P_C$, Lemma~\ref{lem:two-field-calibration} determines
$C$ up to translation, hence determines its Steiner-centered
representative $\overline C=C-s(C)$.
Set
\begin{equation}\label{eq:two-field-geometry-recovery}
                  C_P^0=\overline C+s(P).
\end{equation}
Steiner points add under Minkowski addition, so
$s(P)=s(C)-s(A)$.  Consequently $C_P^0=C-s(A)$, and
\[
 h_{-A_0}=h_P-h_{C_P^0},\qquad
                         \mathcal O_0=\bigcup_P C_P^0.
\]
Any one component gives the same centered footprint.

It remains to recover the probability law, which is not obtained merely
from its support.  Choose any $P=P_C$ and let $v_P=\one_Pv$,
extended by zero to the plane.  The positive component separation
shows that this is the occupation contributed by $C$ alone.
It is bounded, compactly supported, and integrable, and
\begin{equation}\label{eq:two-field-density-convolution}
                        v_P=\one_{C_P^0}*\check\mu_0,
 \qquad \check\mu_0=(z\mapsto-z)_\#\mu_0.
\end{equation}
With Fourier transform $\widehat f(\xi)=\int
e^{-i\xi\cdot x}f(x)\dd x$, one obtains
\begin{equation}\label{eq:two-field-density-quotient}
 \widehat\mu_0(-\xi)
             =\frac{\widehat v_P(\xi)}
                    {\widehat{\one_{C_P^0}}(\xi)}
 \quad\text{whenever }
                    \widehat{\one_{C_P^0}}(\xi)\ne0.
\end{equation}
The denominator is the restriction of an entire function, is
nonzero at the origin, and has a zero set with empty interior in
$\R^2$.  Indeed, vanishing on a real open rectangle would force
vanishing identically by applying the one-variable identity
theorem successively to the two coordinates.  The quotient
therefore specifies the continuous Fourier transform of $\mu_0$
on a dense set and hence everywhere.  Uniqueness of the Fourier
transform of a finite Borel measure recovers $\mu_0$ exactly.
All steps use only the two fields and the stated class bounds.
\end{proof}

\begin{corollary}[Complete fiber, periods, and response completion]
\label{cor:two-field-fiber}
Two admissible configurations and stationary laws have the same
ordered pair $(F_+,F_-)$ almost everywhere if and only if, for some
$h\in\R^2$,
\begin{equation}\label{eq:two-field-fiber}
 \mathcal O'=\mathcal O+h,\qquad A'=A+h,
                         \qquad \mu'=(z\mapsto z+h)_\#\mu.
\end{equation}
These conditions are equivalent to pointwise equality of every
physical forward mean at every nominal center and every finite
command displacement, with the prescribed closed-set bit convention.
Moreover,
\begin{equation}\label{eq:two-field-periods}
                  \operatorname{Per}(F_+,F_-)
                              =\operatorname{Per}(\mathcal O).
\end{equation}
This group is discrete, and it has rank two exactly when
$\mathcal O$ admits a rank-two lattice of periods with finitely
many obstacle orbits.  In \eqref{eq:two-field-periods}, field
periods may be defined almost everywhere or pointwise for the
physical representatives; both definitions give the same group.
Thus the two fixed fields determine both the full forward response
and the full obstacle period group.
\end{corollary}
\begin{proof}
Equality of the two fields gives equality of the canonical
triples by Theorem~\ref{thm:two-field-rigidity} and therefore
\eqref{eq:two-field-fiber}, with $h=s(A')-s(A)$.
Conversely, a common translation couples the actual starts,
segments, and obstacles and preserves every collision bit,
including its solid-start convention.  Recovery of the canonical
triple supplies the defining integral for every other command,
which proves the fiber and completion assertions.

If $w$ is an almost-everywhere common period of the two fields,
the finite prefix inverse shows that $v(\cdot+w)=v$ almost
everywhere.  Translation by $w$ consequently permutes the connected
components $P_C=C+(-A)$ of its measure support.
From $P_C+w=P_{C'}$, cancellation of their support functions gives
$C+w=C'$.  Hence $w$ is a period of $\mathcal O$.
The converse follows directly from translation covariance of the
bit.  Any nonzero period takes a compact obstacle to a different
one, so its length is at least $d$ and the group is discrete.
A discrete subgroup of $\R^2$ has rank at most two; if it has rank
two, local finiteness and bounded obstacle diameters imply that
only finitely many obstacle orbits meet a compact fundamental
cell.  The converse implication is immediate.
\end{proof}

The theorem concerns two exact spatial fields on the plane, and
their field values are still continuum data.  Its data reduction
is nevertheless literal: the command menu consists of the two
vectors $a,-a$, both of a fixed positive length.  The finite
prefix inverse already has a uniform finite-site stability bound.
Quantitative recovery of supports, the negative curvature atoms,
and the Fourier quotient from sampled bits requires the separate
finite estimates below.

\begin{proposition}[A single forward orientation]
\label{prop:two-field-one-orientation}
For any fixed compactly supported launch probability $\mu$, there
are distinct smooth strictly convex single-obstacle configurations
with the same forward response fields in direction $e_1$ at every
positive command length.  Their incoming arcs coincide and their
outgoing arcs differ.
\end{proposition}
\begin{proof}
Choose a nonzero smooth periodic function $b$ whose support is
contained in $\{\phi:\cos\phi>0\}$, and put
\[
 h_\lambda(\phi)=1+\lambda b(\phi),\qquad
               |\lambda|\|b+b''\|_\infty<1/2.
\]
Then $h_\lambda+h_\lambda''>1/2$, so $h_\lambda$ is the support
function of a smooth strictly convex body $C_\lambda$.
On the closed incoming normal semicircle $\cos\phi\le0$, both
$b$ and $b'$ vanish.  The support-point formula
$c_\lambda=h_\lambda n_\phi+h_\lambda' n_\phi^\perp$
therefore gives the same incoming semicircle for every $\lambda$.
Its transverse endpoints are $(0,\pm1)$, and the left endpoint
of every horizontal section at height $w\in[-1,1]$ is
$\ell(w)=-\sqrt{1-w^2}$, including the degenerate endpoint sections.

For any $T>0$, the closed collision convention gives exactly
\[
 B_{Te_1}^{C_\lambda}(u,w)
   =\one_{\{-1\le w\le1,\ \ell(w)-T\le u<\ell(w)\}}.
\]
A free start to the left hits if its segment reaches the left
endpoint; a solid start is zero; and a start to the right of the
section moves away.  The same description includes the two
tangent sections and their solid endpoints.  Thus the bits agree
pointwise for all $T>0$, and integration against the same $\mu$
gives identical response fields at every nominal center.
For $\lambda\ne0$, the support function differs from $h_0$ on
the outgoing normal semicircle.  The bodies are not translates:
a linear translation term vanishing on an open semicircle is zero,
whereas $\lambda b$ is not.  Their outgoing arcs therefore differ.
\end{proof}

\section{Finite geometry from the same two fixed flights}
\label{sec:two-field-finite}
The two-field inverse has a finite realization which keeps both the
flight length and the two flight directions fixed.  Its statistical
step estimates three geometric supports.  Collision locations are not
observed: every point used in a sampled hull is a commanded nominal
center at which the recorded bit was one.  Positivity permits these
same observations to serve all support directions.

Throughout this section $s=6+\beta$, the obstacles have the bounds
\eqref{eq:two-field-geometric-priors}, and the unknown footprint and density obey
\begin{equation}\label{eq:two-field-finite-priors}
 \begin{gathered}
 \diam A\le\Delta<d_0,\qquad
 \|p_A^\circ\|_{C^{6,\beta}}\le M_A,\\
 0<r_A^-\le p_A^\circ+(p_A^\circ)''\le r_A^+,
 \qquad j(z)\ge b_0\dist(z,\partial A)^\gamma
       \quad\hbox{for a.e. }z\in A,
 \qquad \gamma\ge0.
 \end{gathered}
\end{equation}
No density upper bound, continuity modulus, or density evaluations are
supplied.  Fix a dyadic $t>0$ with $2t+\Delta<d_0$, put $a=te_1$,
and use only the two ordinary forward commands $a$ and $-a$.
The unknown law is identical in every attempt.  Write
$P_C=C+(-A)$, $a_0=s(A)$, and $A_0=A-a_0$.
The width of $A$ need not be smaller than an obstacle width.

\subsection{The interface with the technical supplement}
\label{sec:finite-interface}
The exact inverse in Section~\ref{sec:two-field-rigidity} is proved
entirely in this article. For the finite construction we use the
auxiliary estimates in Table~\ref{tab:finite-interface}. We state their
inputs and outputs here to specify the dependency of the argument.
All constants in this subsection depend only on the fixed priors.
The supplement proves the indicated geometric or analytic estimates;
the observation procedures in the present article supply their inputs.

\begin{table}[ht]
\centering
\small
\caption{Auxiliary estimates and their proof locations in the supplement.}
\label{tab:finite-interface}
\begin{tabular}{@{}p{.48\textwidth}p{.46\textwidth}@{}}
\toprule
Estimate used here & Proof in the supplement\\
\midrule
Overlap mass and a boundary layer &
Lemma~\ref{H-lem:cap-mass}\\[3pt]
Complete coarse components and radial brackets &
Lemma~\ref{H-lem:coarse}\\[3pt]
Fixed-accuracy projection normals &
Lemma~\ref{H-lem:rare-coarse-normals}\\[3pt]
Bisection and radial interpolation &
Lemmas~\ref{H-lem:bisection} and \ref{H-lem:radial}\\[3pt]
Signed support smoothing &
Section~\ref{H-sec:stationary}, proof of
Proposition~\ref{H-prop:jitter-modulus}\\[3pt]
Period decisions and common-translation invariance &
Lemmas~\ref{H-lem:recognition}, \ref{H-lem:locking}, and
\ref{H-lem:registration-period-margin}\\
\bottomrule
\end{tabular}
\end{table}

\paragraph{Mass, coarse components and normals.}
Let $P=C+(-A)$, with the smoothness, curvature and density bounds
\eqref{eq:two-field-geometric-priors} and
\eqref{eq:two-field-finite-priors}. There are known $c,e_0>0$ such that
\[
 x\in P_{-e}\ \Longrightarrow\
 \int\one_C(x+z)j(z)\dd z\ge c e^{\gamma+3/2},
 \qquad 0<e<e_0.
\]
The constants are uniform when $A$ varies over its stated class.
For a family of separated convex bodies of bounded diameter and
positive rolling radius, a fixed sufficiently fine grid of labels
which retain every grid node in an inner parallel body and discard every
grid node outside the body gives one approximate hull per complete
component. The construction supplies a center $o$ and fixed
$0<r<R$ with
$\overline B_r(o)\subset P\cap\widehat P\subset\overline B_R(o)$,
and isolated radial intervals enclosing every boundary point.
The aperture has a collar larger than twice the component diameter,
with additional clustering and numerical reserves; discarding clusters meeting the outer cutoff then
keeps every required complete component and no fragment. Refining the
radial data to any prescribed fixed accuracy supplies projection
normals to that accuracy throughout fixed tubular neighborhoods of
the intervals. The interpolation proof permits any fixed normal
tolerance; its displayed value $1/32$ is one choice.
In this article the mass bound and the prefix inverse furnish these
labels, with the threshold and its error reserve specified below.

\paragraph{Bisection and interpolation.}
Suppose a radial interval has width $w_0$ and a label can differ from
the true inside/outside label only within distance $Ae$ of the true
radial endpoint. The label may depend on earlier tests and need not
be monotone. After $k$ midpoint tests the returned radial error is
at most $Ae+2^{-k-1}w_0$. For a body with the common inball,
curvature and $C^{6,\beta}$ support bounds, its radial function has
a uniform $C^{6,\beta}$ bound. From angular spacing $h$ and value
errors at most $e$, the supplement constructs a smooth interpolation
with
\[
 \|\widehat R-R\|_{C^j}\le C(h^{s-j}+eh^{-j}),
 \qquad 0\le j\le3.
\]
When $e\le C_0h^s$, the resulting curve is strictly convex for small
$h$; its Hausdorff error is $Ch^s$ and its support $C^2$ error is
$Ch^{s-2}$. These assertions depend on the stated label and geometric
inputs, irrespective of how the labels were acquired.

\paragraph{Support smoothing.}
There is a fixed signed smooth kernel $\kappa$ on a bounded interval
with integral one and vanishing moments of orders one through six.
For its periodized dilation $T_b$, a function with uniformly bounded
$C^{6,\beta}$ norm satisfies
\[
 \|T_bf-f\|_{C^j}\le Cb^{s-j},\qquad
 \|T_bg\|_{C^j}\le Cb^{-j}\|g\|_\infty,
 \qquad 0\le j\le3.
\]
The supplement gives the explicit dilation coefficients
$8/5,-4/5,8/35,-1/35$ at scales $1,2,3,4$. Taylor's formula and
the moment cancellations give the first estimate, and the $L^1$
norms of the differentiated kernel give the second. The kernel is
used for numerical integration of support functions.

\paragraph{Period locking.}
For the bounded presentation \eqref{eq:two-field-presentation},
put $B=1+L_0$ and $V_*=V_0/r_0$. Complete component estimates with
matched Hausdorff error $e$ in the protected square $Q_{10B}$ suffice
for the following decisions. Under the patch margin
\eqref{eq:two-field-patch-margin}, and below a known positive
threshold for $e$, the tests accept exactly those candidates whose
matched true differences are periods. The accepted true
vectors generate $\Pi$, and $\covol\Pi\ge V_*$. If their lengths
are bounded by the prior-dependent $U_0$, coordinates in an accepted
independent-pair basis have reduced denominators at most
$Q=\max(1,\lceil U_0^2/V_*\rceil)$. Comparison below the rational
gap $Q^{-2}$ recovers the exact relations, the primitive orbit
count, a matched real period basis with error $Ce$, and primitive
free area with error $Ce$. The support comparison and the orbitwise
defect are unchanged by adding one common Minkowski summand or by
one common translation, provided the expanded components remain
separated. Thus the same $\eta$ applies in the canonical frame.
The numerical enclosures, cutoff reserves and periodic assembly
are included in the cited proofs. The finite-configuration result
uses the given complete aperture and omits this period step.

The two-command construction below therefore needs neither the
directional-germ data nor the reciprocal randomized preparations of
the other experiments in the supplement. Each auxiliary hypothesis
has been stated here, and each observation required for the current
finite theorem is supplied by its own two fixed forward channels.


\subsection{Two-direction boundary acquisition}
The prefix inverse in Lemma~\ref{lem:two-field-prefix} already provides
a finite occupation query.  Put
\[
 K=1+\left\lceil (D_0+\Delta)/t\right\rceil.
\]
At a target $x$ its input consists of $2K$ mean values, with nominal
centers on $x+\{0,a,\ldots,Ka\}$.  If every mean has error at most
$\xi$, the computed prefix maximum has error at most $2K\xi$.
Hoeffding's inequality therefore estimates $v(x)$ to error $\eta$
with at most
\begin{equation}\label{eq:two-field-occupation-cost}
 C K^3\eta^{-2}\log\frac{CK}{\varepsilon}
\end{equation}
attempted bits and conditional failure probability at most
$\varepsilon$.  Constants and integer ceilings are understood.
This bound is uniform over the unknown law, and its aperture is fixed.
The coarse threshold uses the known boundary-mass lower bound.
Fix a small prior-dependent depth $d_*>0$ and put
$\lambda=c d_*^{\gamma+3/2}$, using the uniform cap-mass constant
in Lemma~\ref{H-lem:cap-mass}.  Estimate $v$ to error
$\lambda/4$ and retain a grid node when its estimate is at least
$\lambda/2$.  Nodes outside all $P_C$ are discarded, and every
node in $(P_C)_{-d_*}$ is retained.  This is the boundary-layer label
needed in the proof of Lemma~\ref{H-lem:coarse}; its fixed grid,
clustering, interior-ball construction and isolated radial brackets
therefore apply to $P_C$.  The cap-mass estimate is uniform under
\eqref{eq:two-field-finite-priors}, so this construction does not
require knowledge of $A$.

The following refinement uses only the two prescribed directions, even
where they are tangent to the expanded boundary.
Appendix~\ref{sec:two-field-geometric-details} records explicit collar,
strip and finite-grid estimates for the construction, including all
component-assignment and numerical reserves.
\begin{lemma}[A two-direction rare boundary test]
\label{lem:two-field-rare}
There is a fixed collar width $d_c>0$, determined by the stated priors,
with the following property.  Fixed-accuracy occupation data supply
radial brackets for each protected $P_C$ and sufficiently accurate
normals throughout their fixed neighborhoods.  At a target $y$ in such
a neighborhood and an accuracy $0<e<d_c$, a test using only commands
$a,-a$ returns zero surely if $y\notin P_C$, and returns one with
conditional probability at least $1-\varepsilon$ if
$y\in(P_C)_{-e}$.  Its output in the intervening inner layer is
unrestricted.  The test uses at most two nominal centers and
\begin{equation}\label{eq:two-field-rare-cost}
 C e^{-(\gamma+3/2)}\log\frac C\varepsilon
\end{equation}
attempted bits.
\end{lemma}
\begin{proof}
Let $R$ be a fixed upper bound for the curvature radii of $P_C$,
increased if necessary to be positive and larger than $t$.
At a boundary point $y_0$ with outward normal $n$ one has
\begin{equation}\label{eq:two-field-outer-disk}
                 P_C\subset\overline B_R(y_0-Rn).
\end{equation}
Indeed, in support coordinates with the normal angle of $n$ as zero,
\[
 h_{P_C}(\phi)-y_0\cdot n_\phi
   =\int_0^\phi r_{P_C}(u)\sin(\phi-u)\dd u
   \le R(1-\cos\phi),\qquad 0\le\phi\le\pi.
\]
The reverse angular interval proves the same inequality for
$-\pi\le\phi\le0$.  These are exactly the support inequalities
of the disk in \eqref{eq:two-field-outer-disk}.

Choose the fixed collar so that
\[
 d_c<\min\{R/2,t^2/(32R),d_0-\Delta-t\},
\]
and below a fixed interior rolling radius of $P_C$.
Write a collar point uniquely as $y=y_0-dn$, where $|d|\le d_c$.
The coarse construction in Lemma~\ref{H-lem:rare-coarse-normals}, with
its occupation queries replaced by \eqref{eq:two-field-occupation-cost},
can give a rational $\widehat n$ satisfying
$|\widehat n-n|\le\eta_n$ uniformly on each bracket neighborhood,
where a fixed rational $\eta_n>0$ is chosen with
$6\eta_n\le t/(8R)$.  That proof uses only the uniform support,
curvature, and inball bounds of $P_C$: its finite radial interpolation
may be refined to any fixed accuracy before the fine experiment.

Set $E_1=\{e_1,-e_1\}$ and
\[
 V(\widehat n)=\{v\in E_1:\widehat n\cdot v
                  \ge\max_{w\in E_1}\widehat n\cdot w-4\eta_n\}.
\]
Every true maximizer of $n\cdot v$ is in this set.  Every candidate
satisfies $n\cdot v\ge-6\eta_n\ge-t/(8R)$.
For each candidate make ordinary forward attempts from nominal center
$y+tv$ with displacement $-tv$.  Mark the candidate if at least one
bit is one, and return the conjunction of the candidate marks.

Suppose $0<d\le d_c$.  For every candidate,
\begin{align*}
 |y+tv-(y_0-Rn)|^2-R^2
 &=t^2+2t(R-d)n\cdot v-2Rd+d^2\\
 &\ge t^2-t^2/4-t^2/16>0.
\end{align*}
Thus $y+tv$ is outside $P_C$ by
\eqref{eq:two-field-outer-disk}, so every physical start
$y+tv+z$, $z\in A$, is outside $C$.  On the event $y+Z\in C$
the attempted flight ends in $C$ and its bit is one.  Its probability
is at least $c e^{\gamma+3/2}$ if $d\ge e$, by
Lemma~\ref{H-lem:cap-mass}.

Other obstacles cannot occur in these flights.  At $y_0=c_0-z_0$,
the points $c_0\in C$ and $z_0\in A$ have normals $n$ and $-n$.
Every point of a candidate flight has distance at most
$\Delta+t+d_c<d_0$ from $c_0$.
If $d<0$, choose a true maximizing $v_*$.  Since
$n\cdot v_*\ge0$ and $n\cdot(z-z_0)\ge0$ for every $z\in A$,
\[
 n\cdot(y+tv_*+z-u tv_*)
   \ge h_C(n)-d+t(1-u)n\cdot v_*>h_C(n),\quad 0\le u\le1.
\]
This candidate always misses.  Its mark, and hence the conjunction,
is zero surely.  For inner depth at least $e$, a batch of
$C e^{-(\gamma+3/2)}\log(C/\varepsilon)$ attempts at each of the
at most two centers marks both candidates except with probability
$\varepsilon$.  Fresh attempts make the same assertion conditional
on all prior data.  No actual start or direction realization has
been observed.
\end{proof}

Apply this test in the relaxed radial bisection of
Lemma~\ref{H-lem:bisection}.  With radial angular spacing
$h\asymp\epsilon^{1/s}$ and inner layer $e\asymp\epsilon$, the
interpolation of Lemma~\ref{H-lem:radial} recovers all complete protected
expanded bodies with
\begin{equation}\label{eq:two-field-expanded-errors}
 \|\widehat h_{P_C}-h_{P_C}\|_\infty\le C\epsilon,
 \qquad
 \|\widehat h_{P_C}-h_{P_C}\|_{C^2}
                 \le C\epsilon^{(s-2)/s}.
\end{equation}
The attempted-bit cost, including coarse acquisition, is
\begin{equation}\label{eq:two-field-expanded-cost}
 C\epsilon^{-(\gamma+3/2+1/s)}
       \log\frac C\epsilon\log\frac C{\epsilon\delta}.
\end{equation}
These steps construct the expanded bodies before footprint subtraction.
Their proof and conditional confidence allocation are exactly the
radial construction just cited; its label premise is supplied by
Lemma~\ref{lem:two-field-rare}.  Rounding a target to a dyadic grid
of mesh $O(\epsilon)$ changes only the indeterminate geometric layer.
The fixed opposite vectors are exact.

\subsection{Convex hulls of positive nominal records}
Write $S_C^\pm=E_{\pm,C}$ for the closed strips in
\eqref{eq:two-field-strips}, and put
\[
                       Q_C^\pm=\operatorname{conv}(S_C^\pm)+(-A).
\]
Their support functions are those in
Lemma~\ref{lem:two-field-supports}.  The positive response itself may
have a nonconvex support.  The following sampling statement concerns
its convex hull and does not use convexity of that positive set.

\begin{lemma}[Uniform caps and shared hull acquisition]
\label{lem:two-field-hulls}
For each complete protected obstacle and each sign, the convex body
$Q_C^\pm$ can be estimated from the same two fixed commands by an
inscribed rational polygon $\widehat Q_C^\pm$ with Hausdorff error
at most $C\epsilon$, simultaneously with probability at least
$1-\delta$, using
\begin{equation}\label{eq:two-field-hull-cost}
 N\le C\epsilon^{-(\gamma+9/2)}
                         \log\frac C{\epsilon\delta}
\end{equation}
attempted bits.  A dyadic nominal grid of mesh comparable to a
sufficiently small multiple of $\epsilon^{\gamma+9/2}$ suffices.
The same positive records are used for all support directions.
\end{lemma}
\begin{proof}
First fix one strip $S=S_C^+$, the other sign being identical after
reversing $e_1$.  Its parameterization is
\[
 y=c_C(\phi)-u e_1,\qquad
          n_\phi\cdot e_1\le0,\quad 0\le u\le t.
\]
On the incoming open half-circle this parameterization is one-to-one
up to null boundary sets and has Jacobian
$r_C(\phi)|n_\phi\cdot e_1|$.  These assertions also follow by
using the transverse coordinate of each incoming boundary point.
There are uniform $c,r_0>0$ such that
\begin{equation}\label{eq:two-field-strip-ball}
             \area(S\cap B(y,r))\ge c r^3,
                     \qquad y\in S,\quad 0<r<r_0.
\end{equation}
To check the endpoint case as well, represent $y$ by $(\phi_0,u_0)$.
Choose an interval of $u$ of length $c r$ in $[0,t]$ within distance
$Ccr$ of $u_0$.  Choose an angular interval of length $c r$ within
the incoming half-circle and distance $Ccr$ of $\phi_0$, displaced
inward if $\phi_0$ is an endpoint.  It can be chosen so that
$|n_\phi\cdot e_1|\ge c' r$ on this interval.  The upper curvature
radius bounds $|c_C(\phi)-c_C(\phi_0)|$ by $Ccr$.
For a sufficiently small fixed $c$, this parameter rectangle maps
inside $B(y,r)$.  Its Jacobian, the lower curvature radius, and its
two side lengths give \eqref{eq:two-field-strip-ball}.  Take $r_0$
smaller than a fixed multiple of $t$ and of the geometric radii.

For any unit vector $u$, choose a maximizer $y_u\in S$ of
$u\cdot y$.  Equation \eqref{eq:two-field-strip-ball}, with
$r$ a small multiple of $\epsilon$, gives a subset of the strip of
area at least $c\epsilon^3$ whose support deficit in direction $u$
is at most $\epsilon/3$.  The footprint cap satisfies
\begin{equation}\label{eq:two-field-footprint-cap}
 \int_{\{z\in A: h_A(-u)+u\cdot z\le\epsilon/3\}}j(z)\dd z
                       \ge c\epsilon^{\gamma+3/2}.
\end{equation}
Indeed an interior tangent disk of fixed radius in $A$ contains a
rectangle of tangential width $c\sqrt\epsilon$ and normal depth
between $c_1\epsilon$ and $c_2\epsilon$ beneath its contact point.
Every point of a slightly smaller rectangle has distance at least
$c_3\epsilon$ from $\partial A$.  Integration of the stated lower
density bound on this rectangle proves
\eqref{eq:two-field-footprint-cap}.

Let $W$ be a fixed rational rectangle containing the strip differences
needed for all protected components and a collar.  Uniformly choosing
a continuous nominal center $X$ from $W$ and making the positive
command gives
\begin{align}
 &\Prob\{X+Z\in S_C^+,\ X\in\operatorname{cap}(Q_C^+,u,\epsilon)\}
       \nonumber\\
 &\quad=|W|^{-1}\int_Aj(z)
       \int_{S_C^+}\one_{\operatorname{cap}(Q_C^+,u,\epsilon)}(y-z)
                                         \dd y\dd z
       \ge c\epsilon^{\gamma+9/2}.
       \label{eq:two-field-positive-cap}
\end{align}
The cap consists of points $x\in Q_C^+$ satisfying
$u\cdot x\ge h_{Q_C^+}(u)-\epsilon$.  The event displayed uses
only the indicated strip, and every such event gives a positive bit
belonging to $C$, up to null boundary events.  Its obstacle label is
used only in this probability calculation.  The finite assignment
rule below recovers that label from coarse expanded-body data.
The two cap deficits add under Minkowski addition, so
the product of the preceding strip and footprint subsets proves the
inequality.

Replace $X$ by the centers of a uniform finite square grid in $W$.
Conditioning first on $Z=z$, the integrand is the indicator of a
translated strip intersected with a half-plane and $W$.  Each strip
is the difference of a convex swept body and a convex body.  The
perimeters of these bodies, and the lengths of the added straight
cuts in $W$, have a common bound.  Grid cells meeting those
boundaries have total area at most $C\ell$ for mesh $\ell$.
All other cells give the exact indicator integral.  This error bound
is uniform in $z$; integration against the probability density uses
only its total mass.  Consequently the cap probability on the finite
grid differs by at most $C\ell$ from its continuous value.
Choose a dyadic $\ell$ comparable to a sufficiently small multiple
of $\epsilon^{\gamma+9/2}$.  The lower bound in
\eqref{eq:two-field-positive-cap} then remains valid with a smaller
constant.  No regularity modulus of $j$ has been used.

The fixed coarse expanded-body hulls assign records to obstacles.
More precisely choose their Hausdorff error $\sigma_c>0$ so that
$2t+4\sigma_c<d_0-\Delta$.  A positive record of the plus or minus
command belonging to $C$ lies within distance $t+\sigma_c$ of its
coarse hull $H_C$.  Every record belonging to another obstacle has
distance at least $d_0-\Delta-t-\sigma_c$ from $H_C$.
Testing distance at the threshold $t+2\sigma_c$ therefore gives a
unique correct assignment, with a fixed reserve for finite arithmetic.
The surrounding aperture and complete-component cutoff retain every
record from each required complete component.  Exterior fragments are
not used as complete components.  Alternatively, it suffices to retain
one chosen complete component for the calibration below.

For each sign take the convex hull of the assigned positive nominal
centers.  It lies in $Q_C^\pm$ surely.  A fixed-aperture support
function is uniformly Lipschitz in direction, so a net of at most
$C\epsilon^{-1}$ unit directions with spacing $c\epsilon$ suffices
to control its supremum error by $C\epsilon$.  For each net cap the
probability of receiving no assigned positive center after $N$ attempts
is at most $\exp(-cN\epsilon^{\gamma+9/2})$.  The union bound over
the finitely many components, two signs, and this net proves
\eqref{eq:two-field-hull-cost}.  The net is only an analysis device:
the same collection of known nominal centers supplies every support
value of the resulting polygon.  All solid starts and free misses
are counted in $N$ and produce no retained point.
\end{proof}

\subsection{Stable extraction of the contact chord}
The next estimate gives the quantitative step from three support
functions to smooth geometry.  It uses the isolated negative atoms in
\eqref{eq:two-field-curvature-measure}; it does not differentiate an
empirical support function pointwise.

\begin{lemma}[A stable chord correction]
\label{lem:two-field-stable-chord}
Let $C$ satisfy the fixed smoothness and curvature priors.  Let $L$
join its two contact points with normals $e_2,-e_2$, let
$\ell$ be the length of this chord, and let $\theta$ be its normal
angle chosen with $n_\theta\cdot e_1>0$.  Put
$H=h_C-h_L$.  If a computable function $\widehat H$ has
$\|\widehat H-H\|_\infty\le\epsilon$, then finite deterministic
operations produce estimates $\widehat\theta,\widehat\ell$ and a
smooth strictly convex centered body $\widehat C^\circ$ such that
\begin{align}
 |\widehat\theta-\theta|&\le C\epsilon,
             \label{eq:two-field-chord-angle-error}\\
 |\widehat\ell-\ell|&\le C\epsilon^{(s-1)/s},
             \label{eq:two-field-chord-length-error}\\
 \|p_{\widehat C^\circ}-p_C^\circ\|_{C^2}
             &\le C\epsilon^{(s-2)/s},
             \label{eq:two-field-chord-C2-error}\\
 \|p_{\widehat C^\circ}-p_C^\circ\|_\infty
             &\le C\epsilon^{(s-1)/s}.
             \label{eq:two-field-chord-C0-error}
\end{align}
All constants are uniform in the geometric class.  Numerical
integration and representation errors can be included in these bounds.
\end{lemma}
\begin{proof}
The uniform interior radius gives a known lower chord bound
$w_*>0$, while $\ell\le D_0$ and
$n_\theta\cdot e_1\ge w_*/D_0$.  The distributional identity is
\[
 (\partial_\phi^2+1)H
   =r_C(\phi)\dd\phi-\ell(\delta_\theta+\delta_{\theta+\pi}).
\]
For $0<b<\pi/4$ set
\[
 \mathcal D_b f(\alpha)
     =f(\alpha+b)+f(\alpha-b)-2\cos b\,f(\alpha).
\]
Solving the scalar differential equation on $[\alpha-b,\alpha+b]$,
or integrating twice by parts, gives
\begin{equation}\label{eq:two-field-sine-difference}
 \mathcal D_b H(\alpha)
 =\int_{-b}^b\sin(b-|u|)r_C(\alpha+u)\dd u
    -\ell\sum_{\zeta\in\{\theta,\theta+\pi\}:|\zeta-\alpha|<b}
                                   \sin(b-|\zeta-\alpha|).
\end{equation}
If the interval contains no atom this is nonnegative.  If an atom
lies within $b/16$ of its center, it is at most
$-c\ell b+Cb^2$.  Replacing $H$ by $\widehat H$ changes the
expression by at most $4\epsilon$.

Take $b_0=L\epsilon$ with a sufficiently large fixed $L$ depending
on $w_*$.  Evaluate $\mathcal D_{b_0}\widehat H$ on an angular grid
of spacing at most $b_0/8$, with numerical error at most $\epsilon$,
and retain centers where it is smaller than $-w_*b_0/4$.
For sufficiently small $\epsilon$, there is a retained center near
each atom, and every retained center is within $b_0$ of an atom.
The criterion $n_\alpha\cdot e_1\ge w_*/(2D_0)$ selects exactly
the neighborhood of $\theta$, once $\epsilon$ is small.
Choosing one such center proves
\eqref{eq:two-field-chord-angle-error}.
The length of this computational grid has no observation cost.

Fix an even smooth function $\psi$ supported in $(-1,1)$, equal to
one on $[-1/8,1/8]$, and satisfying
$\int u^k\psi(u)\dd u=0$ for $0\le k\le4$.
Such a signed function is obtained from an even central plateau by
adding three pairs of disjoint smooth bumps outside that plateau.
Choose their positive centers distinct.  The matrix of their zeroth,
second and fourth moments is invertible when the supports are small,
by continuity from its Vandermonde limit; choose one such fixed set
and solve the three linear equations.  Odd moments vanish by symmetry.
This construction is fixed and computable.

Put $b=\epsilon^{1/s}$, up to a fixed dyadic factor, and
$\psi_b(\phi)=\psi((\phi-\widehat\theta)/b)$.
For small $\epsilon$ the true atom is in its plateau and the other
atom is outside its support.  Define
\[
 \widetilde\ell
       =-\int\widehat H(\phi)(\psi_b''(\phi)+\psi_b(\phi))\dd\phi.
\]
The true integral equals
$\ell-\int r_C\psi_b$.  Since $r_C\in C^{4,\beta}$ uniformly,
Taylor expansion about $\widehat\theta$ and the vanishing moments
bound the latter integral by $Cb^{5+\beta}=Cb^{s-1}$.
The data error is at most
$C\epsilon(b^{-1}+b)$.  Thus projection of
$\widetilde\ell$ onto $[w_*,D_0]$ proves
\eqref{eq:two-field-chord-length-error}.
An integration error of $O(\epsilon/b)$ is harmless.

Let $L^\circ$ be the centered chord, so
$h_{L^\circ}(\phi)=(\ell/2)|\sin(\phi-\theta)|$.
Let $\widehat L^\circ$ have the estimated parameters and let $T_b$
be convolution with the fixed signed approximation kernel of
Section~\ref{H-sec:stationary}, with moments through degree six
vanishing, at scale $b$.  The function
\[
              \widetilde p=T_b(\widehat H+h_{\widehat L^\circ})
\]
approximates $h_C-m\cdot n_\phi$, where $m$ is the chord midpoint.
For $f_\theta(\phi)=|\sin(\phi-\theta)|$ the curvature identity
$f_\theta''+f_\theta=2(\delta_\theta+\delta_{\theta+\pi})$
gives
\[
 \|T_b f_\theta\|_{C^2}\le Cb^{-1},\qquad
 \|\partial_\theta T_b f_\theta\|_{C^2}\le Cb^{-2}.
\]
The first derivative and supremum norms in the first bound are
bounded independently of $b$.  Consequently the error caused by
estimating the chord is at most
\[
 C\{\epsilon^{(s-1)/s}b^{-1}+\epsilon b^{-2}\}
       =C\epsilon^{(s-2)/s}
\]
in $C^2$.  Smoothing the data error costs $C\epsilon b^{-2}$,
and smoothing the true $C^{6,\beta}$ support costs $Cb^{s-2}$.
The corresponding supremum error is at most $Cb^{s-1}$.
Steiner centering, a bounded linear operation on support functions,
proves \eqref{eq:two-field-chord-C2-error} and
\eqref{eq:two-field-chord-C0-error}.

For small $\epsilon$ the uniform lower curvature radius makes this
smooth centered function an actual strictly convex support.
Its third derivatives remain uniformly bounded: the three error
terms above at derivative order three are bounded by
$C(b^{s-1}b^{-2}+\epsilon b^{-3}+b^{s-3})$.
Finite evaluation and quadratic local interpolation with the fixed
smooth partition used in the retained support construction therefore
give a finite smooth representation with any reserved $O(b^{s-2})$
$C^2$ error and $O(b^{s-1})$ supremum error.  The polygonal input
supports, their known Lipschitz bounds, and the fixed smooth kernels
permit certified finite integration at the precisions just specified.
This completes all numerical operations in the assertion.
\end{proof}

\begin{theorem}[Finite reconstruction from two fixed opposite fields]
\label{thm:two-field-finite-geometry}
Fix the priors \eqref{eq:two-field-finite-priors}, the bounded periodic
class $\mathcal B_\eta$, and one dyadic $t>0$ satisfying
$2t+\Delta<d_0$.  A finite controller using only the two forward
commands $te_1,-te_1$ at one unknown stationary launch law reconstructs
$A_0$ and a smooth strictly convex periodic table approximating
$\mathcal O-a_0$.  Its footprint support error in $C^2$ and its
geometric loss are at most $C\nu$, and its primitive discrete data
are correct, with probability at least $1-\delta$.

For sufficiently small $\nu>0$ and $0<\delta<1/4$, put
\begin{equation}\label{eq:two-field-finite-exponent}
                    Q_{\rm pair}=\frac{(\gamma+9/2)s}{s-2}.
\end{equation}
The attempted-bit and command counts satisfy
\begin{equation}\label{eq:two-field-finite-cost}
 N_\nu\le C\nu^{-Q_{\rm pair}}
              \log\frac C\nu\log\frac C{\nu\delta},
                     \qquad S_\nu\le J_\nu\le N_\nu.
\end{equation}
Every occurrence of a selected nominal center and its sign is included
in $J_\nu$, and every fresh preparation is included in $N_\nu$.
The two flight vectors and their nonzero length remain fixed as
$\nu\downarrow0$.  A nominal dyadic grid of mesh comparable to
\begin{equation}\label{eq:two-field-finite-mesh}
                         \nu^{Q_{\rm pair}}
\end{equation}
suffices.  The controller description is at most
$CJ_\nu\log(C/(\nu\delta))$ bits for fixed computable priors.
The theorem also retains finite approximations to the expanded bodies
$P_C$, their locations, and one pair of raw convex hulls
$\widehat Q_C^+,\widehat Q_C^-$.  In particular a rational polygon
approximating any selected complete $P_C$ to Hausdorff error $C\nu$
can be included in the output without additional observations.

The nominal aperture depends only on the fixed geometric priors.
Actual starts lie in its Minkowski sum with $A$, and flights in the
$t$-parallel body of that sum.  A coarse bound on $|a_0|$ gives an
absolute physical aperture; it does not supply $a_0$.  Computational
support fitting, numerical integration, apparatus calibration,
physical travel, positioning, and realization of the coordinate grid
are not included in the attempted-bit bound. The finite smooth support
representation and the retained rational polygons are the stated
outputs; their arithmetic work and storage are separate resources.
No sharpness of
$Q_{\rm pair}$ is asserted.
\end{theorem}
\begin{proof}
Choose $\epsilon$ comparable to a sufficiently small multiple of
$\nu^{s/(s-2)}$.  First obtain the fixed-accuracy expanded-body
hulls and all complete protected components.  The bounded
presentation supplies a complete obstacle in a prior-bounded nominal
region after any common translation.  Choose one such component $C_*$
for calibration, using the finite coarse hulls and the
complete-component cutoff.  Neither its label nor a free point was
an input.  Enlarge the fixed aperture to contain all required
$P_C$, both strip hulls for $C_*$, their assignment collars, and
all prefix-query sites.  The finite-cloud alternative below supplies
the corresponding complete aperture directly.

Recover the smooth expanded bodies by
\eqref{eq:two-field-expanded-errors}, and the two strip hulls of
$C_*$ by Lemma~\ref{lem:two-field-hulls}.  On their common accuracy
event form
\[
 \widehat H
    =2\widehat h_{P_{C_*}}
       -h_{\widehat Q_{C_*}^+}-h_{\widehat Q_{C_*}^-}
                                  +t|n_\phi\cdot e_1|.
\]
The support identities give
$\|\widehat H-(h_{C_*}-h_{L_*})\|_\infty\le C\epsilon$.
Lemma~\ref{lem:two-field-stable-chord} recovers the centered shape
$C_*^\circ$ with $C^2$ error $C\nu$.  Center the expanded support
and set
\[
  \widehat p_{-A_0}
        =\widehat p_{P_{C_*}}^\circ-
                                    p_{\widehat C_*^\circ}.
\]
This estimates $p_{-A_0}$ in $C^2$ error $C\nu$, since Steiner
centering commutes with Minkowski addition.  Its curvature is positive
for small $\nu$ by the footprint radius reserve.  Reflection gives
the required centered footprint.  Every original component is now
recovered by the same support subtraction
\begin{equation}\label{eq:two-field-finite-subtraction}
       \widehat h_{C-a_0}
                   =\widehat h_{P_C}-\widehat p_{-A_0}.
\end{equation}
The uniform obstacle curvature margin again gives smooth strictly
convex outputs.  Finite smooth representation, with its numerical
errors reserved, is supplied as in the last paragraph of
Lemma~\ref{lem:two-field-stable-chord}.
The observation error in the footprint and all reconstructed
components has support supremum order
$\epsilon^{(s-1)/s}$ and $C^2$ order
$\epsilon^{(s-2)/s}$.  The expanded components themselves retain the
stronger supremum order $\epsilon$.

Apply the retained period-locking and periodic assembly argument to
these original component estimates, in the frame $\mathcal O-a_0$.
The nonperiod-patch margin is invariant under a common translation,
as in Lemma~\ref{H-lem:registration-period-margin}.  The bounds just
obtained tend to zero uniformly, so the fixed locking thresholds are
met.  Primitive relations and orbit counts are exact on the common
event; a matched period basis and free area have error
$O(\epsilon^{(s-1)/s})$, which is $O(\nu)$.

The hull cost \eqref{eq:two-field-hull-cost} dominates the expanded
boundary cost \eqref{eq:two-field-expanded-cost}, since
$\gamma+9/2>\gamma+3/2+1/s$.  Splitting the confidence budget among
coarse acquisition, adaptive boundary queries, and the two hull
experiments proves \eqref{eq:two-field-finite-cost} after substituting
$\epsilon\asymp\nu^{s/(s-2)}$.
The nominal grid in Lemma~\ref{lem:two-field-hulls} is finer than the
$O(\epsilon)$ geometric rounding required in the boundary queries.
Choose it dyadically so that $t$ is an integer multiple of its mesh.
All shifts then preserve the grid, and
\eqref{eq:two-field-finite-mesh} follows.  Only two displacements
occur.  Every sampled grid point has a finite index of
$O(\log(C/\nu))$ bits in the fixed aperture.  Recording these indices,
signs, and batch counts gives the stated description bound.

Finally, the expanded-body estimates have uniformly bounded curvature
radii.  Sampling their smooth boundaries at sufficiently fine
numerical angular spacing and rounding those values gives inscribed
or nearby rational polygons of any reserved $O(\nu)$ Hausdorff error.
This is a finite computation on the already acquired functions and
uses no new collision records.  Their locations are those of the
measured nominal-coordinate bodies $P_C$.
\end{proof}

\begin{corollary}[A complete finite configuration]
\label{cor:two-field-finite-cloud}
Suppose instead there are between one and $n_0$ obstacles, with the
same uniform priors, and a known protected aperture contains every
expanded body $C-A$ and the fixed collars above.  The same controller
recovers $A_0$ and all of $\mathcal O-a_0$ in $C^2$ error $C\nu$
with probability at least $1-\delta$ and the bounds
\eqref{eq:two-field-finite-cost}--\eqref{eq:two-field-finite-mesh}.
No periodicity or positive nonperiod-patch margin is used.
\end{corollary}
\begin{proof}
The complete aperture supplies all coarse expanded components and a
calibration component.  The hull acquisition, chord correction, and
common support subtraction in the theorem recover every body.
Omit the period test and periodic assembly; all preceding estimates
and observation counts remain unchanged.
\end{proof}

\section{Finite recovery of the launch law and prediction}
\label{sec:two-field-finite-law}
The finite experiment also determines the unknown launch law in a
probability metric.  The additional observations use exactly the two
commands of the geometric experiment.  The argument operates on an
isolated occupation component: its normalization is the law of the
sum of a uniform point in the recovered obstacle and the reflected
launch displacement.  Finite moments separate these two factors.
Their conditioning is stated explicitly below.
Appendix~\ref{sec:moment-factor-details} supplies a rational polygon
factor with consistent exact moments, a padded positive programme,
and the full propagation of geometric, sampling and fitting errors.

Write $\mu_0$ for the canonically translated launch probability
measure, and let $W_1$ denote transportation distance for Euclidean
cost.  The moment argument applies to arbitrary compactly supported
probability measures.  The quantitative geometric hypotheses enter
only through the preceding geometric reconstruction.  All windows
and constants in this section are fixed by those hypotheses.

\subsection{Finite occupation moments}
The scalar $a$ in the next two lemmas is an error tolerance; the
command displacement remains the fixed vector $te_1$.
Fix one complete expanded component $P=C_0+(-A_0)$ in the protected
patch.  Its occupation contribution is
\begin{equation}\label{eq:two-field-component-occupation}
 v_C(x)=\int\one_{C_0}(x+z)\,\dd\mu_0(z),\qquad
 \int_{\R^2}v_C(x)\dd x=|C_0|.
\end{equation}
Consequently $v_C/|C_0|$ is the density of $X-Z$, where $X$ is uniform
on $C_0$, $Z$ has law $\mu_0$, and the variables are independent.
This density exists even when $\mu_0$ is singular.

Choose a fixed rational $R\ge1$ sufficiently large that $C_0$, $A_0$,
$P$, and all component cutoffs below lie in $(-R/2,R/2)^2$.
For a multi-index $\alpha=(\alpha_1,\alpha_2)$ put
$x_R^\alpha=(x_1/R)^{\alpha_1}(x_2/R)^{\alpha_2}$.
A geometric estimate with sufficiently small support error supplies
an explicit Lipschitz cutoff $H$ which is one on $P$ and zero on
all other expanded components.  Its Lipschitz constant and support
are uniformly bounded.  For example, apply a piecewise linear
cutoff to the $\ell^\infty$ distance from an estimated polygon for
$P$, with transition width a fixed small fraction of the component
gap.  The error reserve ensures these properties for the true
components.  This construction requires only coarse accuracy for
the cutoff, and $H$ has rational values at rational points.

\begin{lemma}[Moments from finitely many fixed-command bits]
\label{lem:two-field-finite-moments}
Let $m\ge2$, $0<a<1/4$, and $0<\delta<1/4$.  Conditional on a
valid cutoff $H$, the moments
\[
 I_\alpha=\int H(x)v(x)x_R^\alpha\dd x
          =\int v_C(x)x_R^\alpha\dd x,
                    \qquad |\alpha|\le4m,
\]
can all be estimated with error at most $a$ and probability at least
$1-\delta$, using at most
\begin{equation}\label{eq:two-field-moment-cost}
 N_{\rm mom}\le C m^2a^{-4}
                    \log\frac{Cm}{a\delta}
\end{equation}
attempted bits of the two fixed commands.  An integer ceiling is
understood.  Nominal targets form a rational spatial mesh of size
comparable to $a/m$, and the prefix sites extend its fixed window
only by the fixed length $Mt$.  The observations used for the
moments are shared across all multi-indices.
\end{lemma}
\begin{proof}
Put $\ell\asymp a/m$, with a sufficiently small dyadic choice,
and take every mesh node in a fixed square containing the support
of $H$ with a fixed collar.  There are at most $C\ell^{-2}$ nodes.
At each node $x$, estimate the $2M$ means needed by
\eqref{eq:two-field-prefix-inverse}, each to accuracy $c a$.
The maximum of the empirical prefix sums then gives
$|\widehat v(x)-v(x)|\le C M a$; absorb the fixed $M$ by the
choice of the small constant $c$.  Hoeffding's inequality and a
union bound give this event at all mesh nodes using at most
\[
 C\ell^{-2}a^{-2}\log\frac{C}{\ell\delta}
\]
attempted bits.  The rule continues to be valid conditional on
previous observations, because all bits at this stage are fresh.

We verify the spatial quadrature uniformly in the unknown law.
For a fixed $z$, the function being integrated is
$H(x)x_R^\alpha\one_{C_0}(x+z)$.  Its weight has absolute value
at most a constant and Lipschitz constant at most $C(m+1)$.
The total area of mesh cells meeting the boundary of the translated
convex body is at most $C\ell$.  This follows, for example, by
covering the boundary by its inner and outer parallel bands.
On the remaining cells the weight variation is at most $Cm\ell$,
and their total relevant area is bounded.  Thus the Riemann sum
error is at most $Cm\ell$, uniformly for every $z$ in the fixed
footprint bound.  Integration against $\mu_0$ preserves this bound.
Boundary values of an indicator affect only the boundary cells,
so the argument applies also to an atomic launch law.  It uses
neither pointwise continuity of $v$ nor a modulus for $\mu_0$.

Replacing the exact node values by $\widehat v(x)$ adds at most
$Ca$ to every weighted sum.  Decreasing the constants in the
mesh and sampling tolerances proves the accuracy claim and
\eqref{eq:two-field-moment-cost}.  The same empirical node values
are used for all $\alpha$; there is no multiplication of the
observation count by the number of moments.
\end{proof}

\subsection{Quantitative separation of the convolution factors}
The following elementary estimate records the stability needed for
a positive, finite reconstruction of a probability measure.  It also
makes explicit where the nonalgebraic bound in the final acquisition
law enters.  Moment comparison through polynomial approximation is a
classical deconvolution method; see, for example,
\cite[Section~2.3]{RigolletWeed}.  The proof below includes the
bivariate estimate and the perturbation of the recovered convolution
factor required here.  Coordinates in the lemma are dimensionless.

\begin{lemma}[Finite moment separation]
\label{lem:two-field-moment-separation}
Let $U,\widetilde U,Z,\widetilde Z$ be random vectors with probability
laws supported in $[-1,1]^2$, with $U,Z$ independent and
$\widetilde U,\widetilde Z$ independent.  Suppose the moments of
$U$ and $\widetilde U$, and those of $U-Z$ and
$\widetilde U-\widetilde Z$, differ by at most $a$ at every
multi-index of total degree at most $4m$.  Then
\begin{equation}\label{eq:two-field-moment-stability}
 W_1(\mathcal L(Z),\mathcal L(\widetilde Z))
       \le \frac{C}{m}+a(Cm)^{Cm},\qquad m\ge2,
\end{equation}
where $C$ is numerical.  It is enough to assume the bounds on all
the displayed moments; no density or regularity is required.
\end{lemma}
\begin{proof}
For $|\alpha|=k$ the convolution identity is
\begin{equation}\label{eq:two-field-triangular-moments}
 \E(U-Z)^\alpha
   =\sum_{\beta\le\alpha}\binom{\alpha}{\beta}
        \E U^{\alpha-\beta}(-1)^{|\beta|}\E Z^\beta.
\end{equation}
The coefficient of $\E Z^\alpha$ is $(-1)^k$.  Subtract the
corresponding identity for the tilded variables.  All moments of
$U$ and $\widetilde Z$ have absolute value at most one.  If $D_k$
is the largest discrepancy in a moment of $Z$ of total degree $k$,
then $D_0=0$ and
\[
 D_k\le a(1+2^k)+\sum_{r=0}^{k-1}\binom{k}{r}D_r.
\]
Here the binomial identity
$\sum_{|\beta|=r,\,\beta\le\alpha}\binom{\alpha}{\beta}
=\binom{k}{r}$ accounts for both coordinates.  Induction gives
\begin{equation}\label{eq:two-field-factorial-conditioning}
                         D_k\le 2a\,8^k k!.
\end{equation}
Indeed, the sum of the inductive bounds is at most
$2a\,8^k k!(e^{1/8}-1)$, and the forcing term is bounded by the
remaining part of $2a\,8^k k!$.  Thus the triangular inversion has
an explicit finite conditioning bound.

We include the approximation step.  If $f$ is a one-Lipschitz
function on $[-1,1]^2$, subtract its value at zero, so that
$\|f\|_\infty\le2$.  The function
$g(\theta_1,\theta_2)=f(\cos\theta_1,\cos\theta_2)$ is even
in each variable and is Lipschitz on the two-dimensional torus.
Convolve in each variable with the normalized nonnegative kernel
\[
 J_m(\theta)=c_m
       \left(\frac{\sin(m\theta/2)}{\sin(\theta/2)}\right)^4.
\]
Its trigonometric degree is $2(m-1)$, its integral is one, and
\[
 J_m(\theta)\le C\min\{m,m^{-3}|\theta|^{-4}\},\qquad
                 \int_{-\pi}^{\pi}|\theta|J_m(\theta)\dd\theta
                         \le C/m.
\]
The normalization estimate follows by integrating over
$|\theta|\le c/m$ for the lower bound and the displayed elementary
sine bounds for the upper bound.  These bounds also prove the first
moment estimate by splitting at $1/m$.  The tensor convolution
therefore differs from $g$ uniformly by at most $C/m$ and remains
even in each variable.  It has the form
\[
 P(\cos\theta_1,\cos\theta_2),\qquad
 P(x_1,x_2)=\sum_{i,j\le2m}b_{ij}T_i(x_1)T_j(x_2),
\]
where $T_i$ is the Chebyshev polynomial.  Its cosine coefficients
satisfy $|b_{ij}|\le C$.  The recurrence
$T_{i+1}=2xT_i-T_{i-1}$ bounds the sum of the absolute monomial
coefficients of $T_i$ by $3^i$.  Consequently the corresponding
coefficient sum for $P$ is at most $Cm^2 3^{4m}$.

Applying \eqref{eq:two-field-factorial-conditioning} to this
polynomial, whose total degree is at most $4m$, gives
\[
 |\E f(Z)-\E f(\widetilde Z)|
       \le C/m+a(Cm)^{Cm}.
\]
Taking the supremum over Lipschitz functions proves
\eqref{eq:two-field-moment-stability}, by the transportation duality
on a compact set.  For finite laws this is ordinary finite linear
programming duality; approximation on finer finite grids proves the
same identity for arbitrary compact probability measures.
\end{proof}

\subsection{A finite probability output}
Let $N_{\rm geom}(u,\delta)$ denote the observation bound in
Theorem~\ref{thm:two-field-finite-geometry}, with geometric tolerance
$u$.  In particular it supplies a matched body in the canonical
frame, its expanded component, and the protected patch.  The next
statement separates the two costs before substituting that bound.

\begin{theorem}[Finite recovery of the geometry and launch probability]
\label{thm:two-field-finite-law}
Assume the hypotheses of Theorem~\ref{thm:two-field-finite-geometry}.
At one unknown stationary law, the same two fixed opposite commands
admit a finite controller which returns the geometric output of that
theorem and a finitely supported probability measure $\widehat\mu_0$.
For all sufficiently small $\varepsilon>0$ and $0<\delta<1/4$,
with probability at least $1-\delta$ the geometric loss is at most
$C\varepsilon$, the discrete data are correct, and
\begin{equation}\label{eq:two-field-finite-wasserstein}
                         W_1(\widehat\mu_0,\mu_0)
                                      \le C\varepsilon.
\end{equation}
More precisely, let $m=\lceil C_0/\varepsilon\rceil$ and take a
certified dyadic $a_m$ with
\begin{equation}\label{eq:two-field-moment-tolerance}
 2^{-C_1m\log_2(C_1m)-1}
       \le a_m\le 2^{-C_1m\log_2(C_1m)},
\end{equation}
where $C_0,C_1$ are sufficiently large fixed constants.  The
attempted-bit bound is
\begin{equation}\label{eq:two-field-joint-decomposed-cost}
 N\le N_{\rm geom}(c a_m,\delta/2)
       +Cm^2a_m^{-4}\log\frac{Cm}{a_m\delta}.
\end{equation}
Inserting the polynomial geometric bound gives the explicit uniform
sufficient bound
\begin{equation}\label{eq:two-field-joint-exponential-cost}
 N\le\exp\!\left(C\varepsilon^{-1}
                       \log\frac{C}{\varepsilon}\right)
                    \log^2\frac{C}{\delta},
                \qquad S\le J\le N.
\end{equation}
Here every bit, site visit and fixed-command occurrence is charged.
The two displacements remain $a$ and $-a$; their length does not
shrink with $\varepsilon$.  All nominal sites lie in a fixed
prior-dependent aperture.  The coordinates used by the moment stage
have mesh comparable to $a_m/m$; together with the geometric stage
the command description has at most
\begin{equation}\label{eq:two-field-joint-description}
 CJ\left\{\varepsilon^{-1}\log(C/\varepsilon)
                            +\log(C/\delta)\right\}
\end{equation}
bits.  The same coarse bound on the unknown footprint origin as in
the geometric theorem gives a uniform bound on the absolute physical
aperture.  No evaluation of the density, second preparation setting,
angular derivative, or additional displacement vector is supplied.

The probability can be returned with at most
$\binom{4m+2}{2}$ nonzero rational weights at rational grid points;
Corollary~\ref{cor:moment-positive-compression} proves this finite
compression without changing the fitted moments or their error bounds.
The candidate grid has at most $C(m/a_m)^2$ points. The attempted-bit
and command-description bounds do not include deterministic moment
integration, linear-program arithmetic, output storage, apparatus
calibration, physical travel, positioning, or realization of the
required coordinate grid.

The moment stage itself is valid for arbitrary compact launch
probabilities whenever the stated geometric estimates are available.
The exponential bound is a sufficient bound for this two-command
joint inverse; it is not a matched minimax assertion.  Under the
complete finite-configuration hypotheses of
Corollary~\ref{cor:two-field-finite-cloud}, the same acquisition and
probability conclusions hold for all the recovered components,
without a periodicity or nonperiod-patch hypothesis.
\end{theorem}
\begin{proof}
Use half the failure allowance to reconstruct geometry at tolerance
$c a_m$.  Work on its common success event, and take one complete
matched obstacle $\widehat C$ and the rational polygon for its
expanded component.  They give the cutoff $H$ above.  The true body
$C_0$ and $\widehat C$ have Hausdorff distance at most $C c a_m$.
Their areas are uniformly bounded below, and the symmetric
difference of the bodies has area at most $C c a_m$.  The latter
bound follows from the convex parallel-body formula and the uniform
perimeter bound.  Thus the uniform probability measures on the two
bodies have total variation distance at most $C c a_m$.
Every normalized coordinate monomial has absolute value at most one,
so their moments differ by this same bound, independently of degree.
All geometric moments through degree $4m$ can be computed from the
finite geometric representation with certified additional error at
most $c a_m$.  This requires no observations.  For example, uniform
inner and outer rational polygon approximations give area errors and
polynomial integral errors tending to zero with explicit support and
Lipschitz bounds.

Use Lemma~\ref{lem:two-field-finite-moments}, with the remaining
failure allowance, to estimate all $I_\alpha$ to error $c a_m$.
Normalize by a certified positive rational approximation to
$|\widehat C|$, with error at most $c a_m$, to obtain the moments
of $Y=(X-Z)/R$ to error $C c a_m$.  The true moment identities use
$U=X/R$ and $Z/R$, in a fixed dimensionless box; the reconstructed
body supplies the moments of $\widehat U$.  Increase the fixed $R$
if necessary so that all quantities and their reserves are in
$[-1,1]^2$.

We describe an entirely finite positive reconstruction.  Take the
rational grid of spacing $h\asymp c a_m/m$ in $[-1,1]^2$.  It may
be restricted to a known $Ca_m$-parallel neighborhood of the
dimensionless footprint $R^{-1}\widehat A_0$, with its padding larger
than the rescaled certified Hausdorff error plus $2\sqrt2h$. Every
point of $R^{-1}A_0$ then has a nearest retained grid point within $Ch$.  For unknown
nonnegative weights $(p_z)$ summing to one on this finite grid,
form the convolution moments
\[
 T_\alpha(p)=\sum_{\beta\le\alpha}\binom{\alpha}{\beta}
       \widetilde c_{\alpha-\beta}(-1)^{|\beta|}
                         \sum_zp_z z^\beta,
                       \qquad |\alpha|\le4m,
\]
where $\widetilde c$ are the certified moments of $\widehat U$.
Minimize the largest discrepancy between these quantities and the
estimated moments of $Y$.  This is a finite linear program with
rational coefficients; its feasible set is the probability simplex.
A rational solution to any prescribed objective tolerance exists
and can be certified by finite linear programming.

The law of $Z/R$, moved to its nearest retained grid point, is a
comparison probability on this simplex.  Its effect on a monomial
of the difference $\widehat U-Z/R$ of degree at most $4m$ is bounded
by $Cmh$, because all difference coordinates have absolute value
strictly below one after the choice of $R$.  The errors in the
geometric moments contribute at most $2^{4m}c a_m$, if bounded
term by term; this harmless exponential factor can be absorbed
below.  Thus the optimum, with an additional certified tolerance
$a_m$, is at most $C2^{4m}a_m$.
The fitted law and the true law consequently satisfy the assumptions
of Lemma~\ref{lem:two-field-moment-separation} with moment tolerance
at most $C2^{4m}a_m$.  The same factor covers the numerical moment
errors.  The constants in \eqref{eq:two-field-moment-tolerance}
ensure
\[
                  C2^{4m}a_m(Cm)^{Cm}\le 1/m.
\]
After multiplying coordinates by $R$, the fitted law is the claimed
$\widehat\mu_0$, and its transportation error is at most $C/m$.
There are at most $C(m/a_m)^2$ possible output atoms, all with
rational weights and locations.  Numerical construction and the
finite output size are separate from attempted-bit counts.

The union of the two success events has probability at least
$1-\delta$.  Lemma~\ref{lem:two-field-finite-moments} gives
\eqref{eq:two-field-joint-decomposed-cost}.  The polynomial geometric
cost at $c a_m$ and the displayed formula for $a_m$ imply
\eqref{eq:two-field-joint-exponential-cost}; factors polynomial in
$m$ and $\log(1/a_m)$ are absorbed by increasing its fixed constant.
The geometric error $Ca_m$ is smaller than $C\varepsilon$.
Every spatial grid point and its prefix translates use the original
two displacements.  The coordinate denominator has logarithm at
most $C\varepsilon^{-1}\log(C/\varepsilon)$, and the same holds
for the geometric-stage meshes.  Charging repetitions separately
proves the remaining resource and description assertions.  For the
complete finite configuration, the cited corollary provides the same
isolated component and geometric estimates.  Every moment and
probability step above is unchanged and uses no period relation.
\end{proof}

\subsection{Finite prediction of all bounded-length responses}
The probability output has a direct experimental interpretation.
For a bounded nominal window $U$ and a fixed maximum command length
$T$, define the predicted forward family by
\begin{equation}\label{eq:two-field-predicted-family}
 \widehat F_b(x)=\int B_b^{\widehat{\mathcal O}_0}(x+z)
                                      \,\dd\widehat\mu_0(z),
                                  \qquad |b|\le T.
\end{equation}
This is a finite weighted sum for every specified $b$.  The commands
in this predicted family need not have been used in acquisition.

\begin{corollary}[Prediction in local spatial mean]
\label{cor:two-field-finite-prediction}
On the success event of Theorem~\ref{thm:two-field-finite-law},
\begin{equation}\label{eq:two-field-prediction-Lone}
 \sup_{|b|\le T}\int_U
             |\widehat F_b(x)-F_b(x)|\dd x
                         \le C_{U,T}\varepsilon.
\end{equation}
No density upper bound or continuity modulus is required.  More
generally, the right side is $C_{U,T}(\xi+\epsilon)$ if the
matched geometry in the relevant patch has Hausdorff error $\xi$
and $W_1(\widehat\mu_0,\mu_0)\le\epsilon$.  Spatial averages
against any bounded test supported in $U$ inherit this error times
its supremum norm.  The entire predicted family uses the finite
geometric and probability output and needs no further observations.
\end{corollary}
\begin{proof}
Write $D_b(C)=C+[-b,0]$ for the convex swept body.  Up to boundary
sets of planar measure zero, the free-start collision set of a
configuration is
\[
 E_b^{\mathcal O}
   =\left(\bigcup_C D_b(C)\right)\setminus\left(\bigcup_C C\right),
                         \qquad B_b^{\mathcal O}=\one_{E_b^{\mathcal O}}.
\]
Only a uniformly bounded number of bodies meet any fixed relevant
window when $|b|\le T$.  Their swept bodies have perimeters at most
$\operatorname{Per}(C)+2T$.  Hence the collision indicator has a
uniform bound on its total variation on a slightly larger window.
If matched bodies have Hausdorff error at most $\xi$, so do their
swept bodies.  Convex parallel-body area bounds, summed over these
finitely many components, therefore give
\begin{equation}\label{eq:two-field-geometric-indicator-stability}
 \|B_b^{\widehat{\mathcal O}_0}-B_b^{\mathcal O_0}\|_{L^1(V)}
                               \le C_{V,T}\xi.
\end{equation}
For periodic tables, the geometric loss supplies this matching on
every prescribed bounded patch, with a constant depending on that
patch.  No uniform assertion on the whole plane with perturbed
lattice vectors is used.

The translation estimate for a function of bounded variation
(see \cite[Chapter~3]{AFP}) gives
\begin{equation}\label{eq:two-field-BV-translation}
 \int_U |B_b(x+z)-B_b(x+\widetilde z)|\dd x
                              \le C_{U,T}|z-\widetilde z|,
\end{equation}
uniformly for $z,\widetilde z$ in the common bounded footprint box.
For a smooth function this follows by integrating its directional
derivative on segments and applying Fubini.  Approximation in
$L^1$ with convergence of total variations proves the same bound
for the displayed finite-perimeter indicator.  All intermediate
segments stay in a fixed enlarged window, where its variation has
already been bounded.

Integrate the sum of
\eqref{eq:two-field-geometric-indicator-stability} and
\eqref{eq:two-field-BV-translation} against an optimal coupling of
$\mu_0$ and $\widehat\mu_0$, and apply Fubini once more.  This gives
$C_{U,T}(\xi+W_1(\widehat\mu_0,\mu_0))$, uniformly in $b$.
Boundary conventions make no difference to these spatial integrals,
even for an atomic launch law.  The theorem's estimates imply
\eqref{eq:two-field-prediction-Lone}.
\end{proof}

\begin{corollary}[A density estimate under a variation bound]
\label{cor:two-field-finite-density}
Suppose, in addition, that $\mu_0=j_0\dd z$ and that the zero
extension of $j_0$ has total variation at most a known $V$.
There is a finite density output $\widehat j\ge0$, of integral one,
with
\[
                    \|\widehat j-j_0\|_{L^1(\R^2)}
                                    \le C(V+1)\varepsilon
\]
and success probability at least $1-\delta$, using the same two
fixed commands.  A sufficient attempted-bit bound is
\begin{equation}\label{eq:two-field-BV-density-cost}
 \exp\!\left(C\varepsilon^{-2}\log(C/\varepsilon)\right)
                                  \log^2(C/\delta).
\end{equation}
For every fixed bounded nominal window $U$ and $T<\infty$, this
same density output also gives the raw predictions
\begin{align}
 \widehat F_b^{\rm dens}(x)
   &=\int B_b^{\widehat{\mathcal O}_0}(x+z)
                                  \widehat j(z)\dd z,
                       \label{eq:two-field-density-prediction}\\
 \sup_{\substack{x\in U\\ |b|\le T}}
      |\widehat F_b^{\rm dens}(x)-F_b(x)|
   &\le C_{U,T,V}\varepsilon.
                       \label{eq:two-field-uniform-prediction}
\end{align}
The variation hypothesis concerns these strong density and uniform
raw-mean conclusions.  It is not needed for the probability or
local spatial-mean conclusions.
\end{corollary}
\begin{proof}
Apply Theorem~\ref{thm:two-field-finite-law} at transportation
tolerance $c\varepsilon^2$.  Let $\rho$ be a fixed nonnegative,
smooth compactly supported probability kernel, put
$\rho_h(z)=h^{-2}\rho(z/h)$, and output
$\widehat j=\rho_\varepsilon*\widehat\mu_0$.
This is a finite mixture of explicitly specified densities.
The translation estimate for $BV$ functions gives
\[
 \|\rho_h*j_0-j_0\|_1\le CVh.
\]
For any coupling of the two laws,
$\|\rho_h(\cdot-z)-\rho_h(\cdot-\widetilde z)\|_1
\le C h^{-1}|z-\widetilde z|$.
Integrating an optimal coupling therefore gives
\[
 \|\rho_h*\widehat\mu_0-j_0\|_1
          \le C\{Vh+h^{-1}W_1(\widehat\mu_0,\mu_0)\}.
\]
Take $h=\varepsilon$ and substitute the theorem's cost at
$c\varepsilon^2$.

For the uniform raw predictions, first observe the planar $BV$
bound
\[
 \|j_0\|_2
 \le\int_0^\infty |\{j_0>q\}|^{1/2}\dd q
 \le C\int_0^\infty\operatorname{Per}(\{j_0>q\})\dd q
 =C|Dj_0|(\R^2)\le CV.
\]
The first inequality is Minkowski's integral inequality applied to
the layer-cake representation, the second is the planar
isoperimetric inequality, and the equality is the coarea formula.
For fixed $U,T$, all evaluations $x+z$ and their commanded segments
lie in a fixed enlarged canonical window.  Its relevant obstacle
count is bounded by the geometric priors.  The matching statement
and the swept-body estimate
\eqref{eq:two-field-geometric-indicator-stability} therefore imply,
uniformly for $x\in U$ and $|b|\le T$,
\[
 \int j_0(z)
  |B_b^{\widehat{\mathcal O}_0}(x+z)
       -B_b^{\mathcal O_0}(x+z)|\dd z
                 \le C_{U,T}\|j_0\|_2\sqrt\xi
                 \le C_{U,T}V\sqrt\xi,
\]
where $\xi$ is the geometric tolerance.  The bound follows by
Cauchy--Schwarz and the area bound on the indicator's symmetric
difference.  Replacing $j_0$ by $\widehat j$ in the estimated
geometry adds at most $\|\widehat j-j_0\|_1$ to any raw mean.
In the present construction $\xi=Ca_m$ with
$m\asymp\varepsilon^{-2}$, so $\sqrt\xi\le C\varepsilon$.
These bounds prove \eqref{eq:two-field-uniform-prediction}.
The use of a fixed enlarged window also accounts for perturbations
of lattice vectors in a periodic output.
\end{proof}

\appendix
\section{Geometric estimates for finite two-command acquisition}
\label{sec:two-field-geometric-details}
This section gives the uniform geometric estimates used in
Section~\ref{sec:two-field-finite}.  The constants depend only on
\eqref{eq:two-field-geometric-priors},
\eqref{eq:two-field-finite-priors}, and the fixed positive length
$t$.  In particular they do not depend on a modulus of continuity
of the unknown density.  Put
\[
 r_- =\kappa_+^{-1},\qquad r_+=\kappa_-^{-1},\qquad
 g=d_0-\Delta>2t,\qquad D_P=D_0+\Delta.
\]
An expanded component has curvature radius between
$r_-+r_A^-$ and $r_++r_A^+$, diameter at most $D_P$, and
distance at least $g$ from every other expanded component.

\subsection{A fixed collar and the conjunction of two tests}
\begin{lemma}[Quantitative tangential collar]
\label{lem:two-field-collar-details}
Choose $R>\max(t,r_++r_A^+)$, and fix
\[
 0<d_c<\min\left\{\frac{r_-+r_A^-}{4},\frac R2,
                 \frac{t^2}{32R},g-t\right\},
 \qquad 0<\eta_n\le\frac{t}{48R}.
\]
Let $y=y_0-dn$, where $y_0\in\partial P_C$, $n$ is its outward
unit normal, and $|d|\le d_c$.  Suppose that a rational vector
$\widehat n$ satisfies $|\widehat n-n|\le\eta_n$.  Define
\[
 V=\left\{v\in\{e_1,-e_1\}:\widehat n\cdot v
     \ge\max_{w\in\{e_1,-e_1\}}\widehat n\cdot w-4\eta_n\right\}.
\]
For every $v\in V$, use nominal center $y+tv$ and displacement
$-tv$.  If $d<0$, at least one of these experiments has collision
probability zero.  If $e\le d\le d_c$, with $0<e<e_0$ for a fixed
$e_0>0$, every one has collision probability at least
$q_e=c e^{\gamma+3/2}$, where $c>0$ is known from the priors.
The rule which returns one if and only if each of the at most two
batches contains a collision therefore has error probability at most
$\varepsilon$ outside the inner layer of width $e$, using at most
\[
 2\left\lceil q_e^{-1}\log(2/\varepsilon)\right\rceil
\]
attempts.  These statements hold conditionally on successful coarse
normal acquisition and on all subsequent adaptive history.
\end{lemma}
\begin{proof}
The lower rolling radius makes the signed normal coordinates unique
in the displayed collar.  The upper curvature-radius bound gives
\[
             P_C\subset\overline B_R(y_0-Rn).
\]
For clarity, if the normal angle of $n$ is zero, the support of
$P_C-y_0$ equals
$\int_0^\phi r_{P_C}(u)\sin(\phi-u)\dd u$ for
$0\le\phi\le\pi$, and is at most $R(1-\cos\phi)$.
Integrating along the reverse angular interval gives the same bound
for $-\pi\le\phi\le0$.  These are the support inequalities for
the asserted disk.

Every true maximizer $v_*$ of $n\cdot v$ over the two directions
belongs to $V$: replacing $n$ by $\widehat n$ changes a comparison
by at most $2\eta_n$.  For every $v\in V$,
\[
             n\cdot v\ge n\cdot v_*-6\eta_n
                           \ge-\frac{t}{8R}.
\]
Consequently, if $0<d\le d_c$,
\begin{align*}
 |y+tv-(y_0-Rn)|^2-R^2
  &=t^2+2t(R-d)n\cdot v-2Rd+d^2\\
  &\ge t^2-\frac{t^2}{4}-\frac{t^2}{16}
      =\frac{11t^2}{16}>0.
\end{align*}
Thus $y+tv\notin P_C$, including when the command is nearly tangent
to $\partial P_C$.  Since $P_C=C-A$, every physical start
$y+tv+z$, $z\in A$, lies outside the closed obstacle $C$.

Write $y_0=c_0-z_0$ for the support points of $C$ and $A$ with
normals $n$ and $-n$.  Every point of every candidate segment has
the form
\[
          c_0+(z-z_0)-dn+(1-u)tv,
                   \qquad z\in A,\quad 0\le u\le1.
\]
Its distance from $c_0$ is at most $\Delta+d_c+t<d_0$.
Thus no other obstacle meets any candidate segment.  When $d<0$,
the true maximizing candidate has $n\cdot v_*\ge0$, and
\[
 n\cdot\{y+tv_*+z-utv_*\}
 =h_C(n)-d+n\cdot(z-z_0)+(1-u)t n\cdot v_*
 >h_C(n).
\]
It always misses.  This proves the exterior assertion for the
conjunction; individual nonmaximizing candidates need not have zero
collision probability there.

For positive depth, the event $y+Z\in C$ forces a collision from
every free shifted start just established.  The overlap estimate
of Lemma~\ref{H-lem:cap-mass} gives probability at least
$q_e=c e^{\gamma+3/2}$.  Its constant is uniform here: the proof
uses only lower rolling radii, upper diameter and support bounds,
and $b_0,\gamma$.  The footprint $K$ in that lemma may therefore
range over the present class of unknown bodies $A$.  More
explicitly, replacing $A$ by $A_{-e/4}$ leaves a uniform rolling
radius, gives an overlap of area at least $c e^{3/2}$, and on that
overlap the density is at least $b_0(e/4)^\gamma$.

Fresh preparations have the same law at every candidate.  A batch
of the stated size has no collision with probability at most
$(1-q_e)^{\lceil q_e^{-1}\log(2/\varepsilon)\rceil}
\le\varepsilon/2$.  The union bound over $V$ proves the assertion.
The coarse event fixes the validity of the normal estimate; fresh
draws then prove the same bound after conditioning on the entire
previous adaptive record.  For a rounded nominal target the lemma
is applied at that actual target, which must stay in the fixed
normal neighborhood.  The signed depth is a Lipschitz function,
so rounding by $O(e)$ enlarges the intended radial bisection
ambiguity by $O(e)$.  The sure-zero exterior assertion refers to
the actual rounded target.
\end{proof}

\subsection{Area at grazing points of a collision strip}
\begin{lemma}[A uniform strip neighborhood]
\label{lem:two-field-strip-details}
For either sign let $S_C^\pm$ be the closed collision strip of
Section~\ref{sec:two-field-finite}.  There are known constants
$c_S,r_S>0$ such that
\[
        \area(S_C^\pm\cap B(y,r))\ge c_S r^3,
              \qquad y\in S_C^\pm,\quad 0<r<r_S.
\]
This estimate includes both tangent endpoints and both ends of each
flight segment.
\end{lemma}
\begin{proof}
It suffices to consider the positive direction.  In normal angle,
\[
 F(\phi,u)=c_C(\phi)-u e_1,
   \qquad \frac\pi2\le\phi\le\frac{3\pi}2,
   \quad 0\le u\le t.
\]
The transverse coordinate $c_C(\phi)\cdot e_2$ is strictly
decreasing on the open angular interval because its derivative is
$r_C(\phi)\cos\phi<0$.  The parametrization is consequently
injective there.  Its absolute Jacobian is
$r_C(\phi)|\cos\phi|$, and its image is the strip, apart from
boundary curves of area zero.

Fix $y=F(\phi_0,u_0)$, and reduce $r_S$ so that
$r<t$ and
\[
 \alpha=\frac{r}{16(1+r_+)}<\frac{\pi}{12},
       \qquad v_r=\frac r{16}<\frac t4.
\]
Move $\phi_0$, if necessary, to the closest point $\phi_*$ of
$[\pi/2+2\alpha,3\pi/2-2\alpha]$, and put
$I=[\phi_*-\alpha/2,\phi_*+\alpha/2]$.
Then $|I|=\alpha$, $|\phi-\phi_0|\le5\alpha/2$ on $I$,
and $|\cos\phi|\ge\alpha/2$ there.  Similarly move $u_0$
to the closest point $u_*$ of $[v_r,t-v_r]$ and put
$J=[u_*-v_r/2,u_*+v_r/2]$.
Thus $|J|=v_r$ and $|u-u_0|\le3v_r/2$.

Since $|c_C'(\phi)|\le r_+$, every point of $F(I\times J)$
is at distance at most
\[
        \frac52r_+\alpha+\frac32v_r<r
\]
from $y$.  The area formula on this injective rectangle now gives
\[
 \area(S_C^+\cap B(y,r))
 \ge\int_I\int_J r_C(\phi)|\cos\phi|\dd u\dd\phi
 \ge\frac{r_-}{2}\alpha^2v_r=c_Sr^3.
\]
The construction deliberately moves the angular rectangle into the
incoming half-circle.  At a tangent endpoint the Jacobian loses one
power of $r$, which is exactly the third power in the bound.
Reflection in the first coordinate proves the negative case.
\end{proof}

\subsection{Finite nominal grids and complete components}
\begin{lemma}[Boundary cells and positive cap probabilities]
\label{lem:two-field-grid-details}
Let $W$ be a fixed dyadic rectangle containing
$P_C+t\overline B$ for every component to be retained, with a
fixed positive collar.  Put
$\kappa=\gamma+9/2$ and
$Q_C^\pm=\operatorname{conv}(S_C^\pm)-A$.
For every unit direction $u$ and sufficiently small $\epsilon>0$,
a uniform nominal center in $W$, followed by the corresponding
fixed command, has probability at least $c\epsilon^\kappa$ of
producing a positive record from $C$ in the support cap
\[
 \{x\in Q_C^\pm:
       u\cdot x\ge h_{Q_C^\pm}(u)-\epsilon\}.
\]
The same assertion holds, with a smaller $c$, when the nominal center
is uniform among the centers of the square grid cells of side
$\ell$ partitioning $W$, provided
$\ell\le c_0\epsilon^\kappa$.  No bound on the density other
than \eqref{eq:two-field-finite-priors} is used.
\end{lemma}
\begin{proof}
Choose a point $y_u\in S_C^\pm$ maximizing $u\cdot y$.
Lemma~\ref{lem:two-field-strip-details}, applied with radius
$\epsilon/6$, gives strip area at least $c\epsilon^3$ with
support deficit at most $\epsilon/6$.

At the footprint support point $z_0$ with outward normal $-u$,
the interior disk of radius $r_A^-$ lies in $A$.  With $v=-u$,
use coordinates
\[
 z=z_0-qv+w v^\perp,\qquad
 \frac\epsilon8\le q\le\frac\epsilon6,
 \qquad |w|\le\frac{\sqrt{r_A^-\epsilon}}{16}.
\]
For $\epsilon$ sufficiently small relative to $r_A^-$ this
rectangle has distance at least $\epsilon/64$ from the boundary
of the interior disk, hence from $\partial A$.  For example its
squared distance from the disk center is
$(r_A^--q)^2+w^2\le(r_A^--\epsilon/64)^2$ after decreasing
the fixed upper bound for $\epsilon$.
Its area is $c\epsilon^{3/2}$ and its support deficit in
direction $-u$ is $q\le\epsilon/6$.  Its probability mass is
therefore at least $c\epsilon^{\gamma+3/2}$.

For every pair $(y,z)$ in these two subsets, $x=y-z$ belongs to
$W$ and has support deficit at most $\epsilon/3$ in $Q_C^\pm$.
Fubini's theorem, divided by $|W|$, proves the continuous lower
bound $c\epsilon^{\gamma+9/2}$.  Strip boundary points can be
removed here: they have zero area in the integration over $y$.

We give the grid argument with its conditioning order explicit.
Fix $Z=z$.  Membership of a nominal center in the relevant event
is determined by a translated collision strip, one half-plane,
and $W$.  The condition $x\in Q_C^\pm$ adds no boundary, since
$x\in S_C^\pm-z$ and $z\in A$ already imply it.
Up to boundary conventions the strip indicator is the difference
of the indicators of $C+[\mp te_1,0]$ and $C$.
Each of these convex sets has a uniformly bounded perimeter; for
example its perimeter is at most $\pi(D_0+t)$.
The additional half-plane cuts and the sides of $W$ have bounded
total length.

Only grid cells meeting one of these boundaries can disagree with
the value of the indicator at the cell center.  A curve of length
$L$ meets cells of total area at most
$C(L\ell+\ell^2)$: partition the curve by arc length into
at most $1+L/\ell$ subarcs of length at most $\ell$; each
subarc and all cells meeting it are contained in a disk of radius
$3\ell$ about its initial point.  This also covers cells whose
centers lie exactly on a boundary.  Thus the discrepancy between
the uniform continuous probability and the uniform grid probability
is at most $C\ell$, uniformly in the fixed $z$.
Integration against the probability law of $Z$ preserves this
bound.  Taking $\ell\le c_0\epsilon^\kappa$ with $c_0$
small enough preserves half the continuous lower bound.
\end{proof}

\begin{lemma}[Assignment and complete hull acquisition]
\label{lem:two-field-assignment-details}
Suppose that the successful coarse event supplies complete expanded
components $P_C$ with convex hull estimates $H_C$ satisfying
$d_H(H_C,P_C)\le\sigma$, where
\[
            0<\sigma\le\frac{g-2t}{16}.
\]
A positive nominal record is assigned to $C$ if the upper endpoint
of a distance enclosure for $\dist(x,H_C)$ is less than
$t+2\sigma$, using an enclosure of width at most $\sigma/4$.
This rule assigns all
records from $C$ correctly and never assigns another component's
record to $C$.  It remains valid when only the required complete
components are retained and all other records are discarded.

With the grid of Lemma~\ref{lem:two-field-grid-details}, the convex
hulls of the assigned positive nominal centers have simultaneous
Hausdorff error at most $C\epsilon$ after
\[
       C\epsilon^{-\kappa}\log\frac C{\epsilon\delta}
\]
attempts, for the two signs and all of the prior-bounded number of
required complete components, with conditional probability at least
$1-\delta$.
\end{lemma}
\begin{proof}
A positive record $x$ from $C$ belongs to $S_C^\pm-A$, so
$\dist(x,P_C)\le t$.  Hence
\[
 \dist(x,H_C)\le t+\sigma.
\]
If the same record belongs instead to a different component $D$,
then
\[
 \dist(x,H_C)\ge g-t-\sigma
       \ge t+15\sigma.
\]
The threshold has a reserve of $\sigma$ on the first side and
$13\sigma$ on the second.  The specified numerical enclosure
preserves both decisions.  A segment of length $t<d_0$ cannot
meet two distinct obstacles, so every positive record has a unique
physical component.  That component label is established by this
distance test; it is not a returned part of the observation.

Here is a concrete completeness cutoff.  Perform the coarse
occupation search on a larger square $U$, with grid spacing and
depth thresholds fixed from the priors.  On its accuracy event all
retained nodes lie in their assigned $P_C$, all sufficiently deep
nodes are retained, and the clustering and interior-ball argument
of Lemma~\ref{H-lem:coarse} applies.  Discard every cluster hull
whose distance from $\partial U$ is at most $D_P+2\sigma$.
If $P_C$ meets the exterior of $U$, select a point of $P_C$ on
one exterior side.  The diameter bound $D_P$ forces every retained
node of this component, and its entire cluster hull, to lie in the
width-$D_P$ strip along that side of $U$.  Such a fragment is
discarded.  Thus every surviving hull represents a complete
component.  Taking $U$ with a collar exceeding $2D_P+4\sigma$
about the required region protects its complete components from
this cutoff.  Finally enlarge $W$, still by a fixed amount, to
contain their $P_C+t\overline B$ and a positive collar.  All
their possible nominal collision records are then sampled in
$W$.  Fragments of other components cannot pass the distance
test just proved, whether or not those fragments have been
retained in the coarse list.

For each required $C$ and sign, take a direction net of spacing
$c\epsilon$ on the circle, of cardinality at most
$C\epsilon^{-1}$.  The net is used only in the proof.  The
probability that a particular cap receives no assigned positive
record in $N$ independent attempts is at most
$\exp(-cN\epsilon^\kappa)$ by
Lemma~\ref{lem:two-field-grid-details} and the correct assignment.
The union bound over the net, components, and signs gives a record
in every such cap at the asserted $N$.

Every retained center lies in $Q_C^\pm$, so its convex hull is
inscribed.  Both the true and empirical supports have a uniform
Lipschitz constant in the direction, since the corresponding sets
lie in the fixed bounded rectangle $W$.  Extending the cap
inequalities from the net therefore gives supremum support error
$C\epsilon$, equivalently Hausdorff error $C\epsilon$.
All retained vertices are known dyadic nominal centers.  Neither
collision locations nor a sample of the launch displacement is
needed for their rational polygon representation.
\end{proof}

\subsection{Signed plateaux and curvature-atom mass}
\begin{lemma}[A fixed plateau with zero moments]
\label{lem:two-field-plateau-details}
There is an even $C^\infty$ function $\psi$, supported in
$(-1,1)$ and equal to one on $[-1/8,1/8]$, such that
\[
             \int u^k\psi(u)\dd u=0,
                         \qquad 0\le k\le4.
\]
The function can be fixed once by an invertible three-dimensional
linear system.  If
\[
 (\partial_\phi^2+1)H
       =r(\phi)\dd\phi-\ell(\delta_\theta+\delta_{\theta+\pi}),
       \qquad \|r\|_{C^{4,\beta}}\le M,
\]
$\|\widehat H-H\|_\infty\le\epsilon$, and
$|\widehat\theta-\theta|\le C_0\epsilon$, then, with
$b\asymp\epsilon^{1/(6+\beta)}$,
\[
 -\int\widehat H(\phi)
   \left\{b^{-2}\psi''\left(\frac{\phi-\widehat\theta}{b}\right)
             +\psi\left(\frac{\phi-\widehat\theta}{b}\right)\right\}
             \dd\phi
       =\ell+O(\epsilon^{(5+\beta)/(6+\beta)}).
\]
The implied constant depends only on $M,C_0$ and the fixed function
$\psi$.
\end{lemma}
\begin{proof}
Take an even smooth central function $\psi_0$, equal to one on
$[-1/8,1/8]$ and supported in $(-1/4,1/4)$.  Fix an even
nonnegative smooth function $k$, supported in $(-1,1)$ and with
integral one.  Put
\[
 x_1=\frac13,\qquad x_2=\frac12,\qquad x_3=\frac23,
            \qquad 0<\rho<\frac1{64},
\]
and define disjoint even pairs of bumps
\[
 B_j(u)=\rho^{-1}\left\{
       k\left(\frac{u-x_j}{\rho}\right)
          +k\left(\frac{u+x_j}{\rho}\right)\right\}.
\]
Writing $m_i=\int u^i k(u)\dd u$, the column indexed by $j$
of the matrix of moments of orders zero, two and four is
\[
 2\begin{pmatrix}
 1\\ x_j^2+\rho^2m_2\\
 x_j^4+6\rho^2m_2x_j^2+\rho^4m_4
 \end{pmatrix}.
\]
Subtracting suitable multiples of the preceding rows gives its
determinant exactly as
\[
                 8\prod_{i<j}(x_j^2-x_i^2)>0.
\]
Choose the three coefficients of $\sum_j c_jB_j$ to cancel
the corresponding moments of $\psi_0$.  The sum
$\psi=\psi_0+\sum_jc_jB_j$ has the asserted support and plateau;
odd moments vanish by symmetry.  All derivative and integral norms
of this single fixed function are finite.  This is a signed test
function; positivity is neither asserted nor used.

Let $\psi_b(\phi)=\psi((\phi-\widehat\theta)/b)$, interpreted
in the local angular interval containing $\widehat\theta$.
For sufficiently small $\epsilon$, the atom at $\theta$ lies
in the plateau because $C_0\epsilon<b/8$, while the opposite
atom lies outside its support.  Distributional integration by
parts gives
\[
       -\int H(\psi_b''+\psi_b)
             =\ell-\int r\psi_b.
\]
Taylor's formula for $r\in C^{4,\beta}$ about
$\widehat\theta$, with the five vanishing moments, gives
$|\int r\psi_b|\le C M b^{5+\beta}$.
The replacement of $H$ by $\widehat H$ costs at most
\[
 \epsilon\bigl(b^{-1}\|\psi''\|_1+b\|\psi\|_1\bigr).
\]
At the stated scale both leading errors have order
$\epsilon^{(5+\beta)/(6+\beta)}$.  A numerical integration
error bounded by a fixed multiple of $\epsilon/b$ is of this
same order and may be reserved in advance.
\end{proof}

The angular search in Lemma~\ref{lem:two-field-stable-chord}
supplies the $O(\epsilon)$ location error required here.  Its sine
second difference has nonnegative background at every atom-free
interval, whereas an atom within one sixteenth of the interval
radius contributes at most $-c\ell b_0$.  Choosing
$b_0=L\epsilon$ makes that contribution exceed the
$4\epsilon$ data error and the $O(b_0^2)$ smooth background.
The positive lower bound for the chord's vertical projection
separates the two opposite angular neighborhoods uniformly.

For completeness, the smoothing used after mass extraction is the
periodic convolution $T_b$ with
\[
 K_b=\frac85 k_b-\frac45 k_{2b}
                  +\frac8{35}k_{3b}-\frac1{35}k_{4b},
                 \qquad k_b(u)=b^{-1}k(u/b).
\]
Its integral is one and its nonconstant moments through degree six
vanish.  Taylor's formula gives
\[
 \|T_bf-f\|_{C^j}\le Cb^{s-j}\|f\|_{C^{6,\beta}},
   \qquad
 \|T_bq\|_{C^j}\le Cb^{-j}\|q\|_\infty,
                 \qquad 0\le j\le2.
\]
The chord support has two point masses as curvature, so its smoothed
$C^2$ norm is $O(b^{-1})$ and its angular derivative has smoothed
$C^2$ norm $O(b^{-2})$.  Combining these two bounds with the
mass error $O(\epsilon^{(s-1)/s})$ and angular error
$O(\epsilon)$ gives the $C^2$ error
$O(\epsilon^{(s-2)/s})$ in
Lemma~\ref{lem:two-field-stable-chord}.
All integrations here act on acquired finite support representations
and fixed smooth kernels.  They use no additional observations.

\section{Moment reconstruction with an estimated geometric factor}
\label{sec:moment-factor-details}
The finite law theorem uses three different approximations: a geometric
factor is reconstructed, occupation moments are sampled, and a positive
probability is fitted to those moments.  We give quantitative interfaces
between these operations.  In particular, the fitted geometric moments
can be the exact moments of a rational polygon.  They then belong to an
actual probability measure at every degree used in the reconstruction.

Fix the complete component used in
Theorem~\ref{thm:two-field-finite-law}.  The constants in the lemmas below
depend only on a bound for its aperture, a lower bound $g_*>0$ for the
gap between expanded components, a lower bound $v_*>0$ for the obstacle
area, an upper bound $p_*$ for the relevant convex perimeters, and the
fixed prefix length $M$.  These bounds are supplied by the quantitative
geometric priors.  All assertions about estimated geometry are made on
the geometric stage's common accuracy event.  Numerical integration and
rational arithmetic are deterministic operations conditional on that
event.

Enlarge the fixed rational scale $R$ of
Section~\ref{sec:two-field-finite-law}, if necessary, so that the selected
body, the centered footprint, and their small approximation collars lie
in $(-R/4,R/4)^2$.  The cutoff support lies in $(-R/2,R/2)^2$.
This is a choice of coordinates within the fixed aperture.  We shall use
candidate normalized displacements in $[-1/2,1/2]^2$; consequently all
coordinate differences appearing in the factorization have absolute
value less than one.  For a finite signed measure, $\|\cdot\|_{\rm TV}$
denotes its full variation norm.

\subsection{A complete component and a rational geometric factor}
\begin{lemma}[A protected cutoff with rational values]
\label{lem:moment-protected-cutoff}
Let $P=C_0+(-A_0)$ be a complete expanded component, and let
$\widehat P$ be a rational convex polygon with
$d_H(\widehat P,P)\le\xi$.  Choose a fixed positive rational $r$ with
$r<g_*/(16\sqrt2)$ and suppose $\xi<r$.  There is an explicitly
specified cutoff $H$, with $0\le H\le1$, which is one on $P$, zero
on every other expanded component, and has Lipschitz constant at most
$r^{-1}$.  Its support is contained in
$\widehat P+[-2r,2r]^2$, and its values at rational points are rational.
For every compact launch probability with support $A_0$, including a
singular probability, one has the pointwise identity
\begin{equation}\label{eq:moment-cutoff-identity}
 H(x)v(x)=v_C(x)=\int\one_{C_0}(x+z)\,\dd\mu_0(z).
\end{equation}
\end{lemma}
\begin{proof}
Let $d_\infty(x,\widehat P)$ be distance for the maximum norm, and put
\[
 \chi(u)=
 \begin{cases}
 1,&0\le u\le1,\\
 2-u,&1<u<2,\\
 0,&u\ge2,
 \end{cases}
 \qquad
 H(x)=\chi\bigl(d_\infty(x,\widehat P)/r\bigr).
\]
Distance to a closed set is one-Lipschitz.  The bound on $\xi$ gives
$H=1$ on $P$.  If $P'$ is another expanded component, then
\[
 \inf_{x\in P'}d_\infty(x,\widehat P)
       \ge g_*/\sqrt2-\xi>2r,
\]
so $H$ is zero there, including the boundary.  The support and
Lipschitz assertions follow directly from the definition.  Distance
from a rational point to a rational polygon in the maximum norm is
the optimum of a finite rational linear programme; its optimum is
rational.  Thus $H$ has the asserted rational values.

For any component $C'$, its occupation contribution vanishes outside
$C'+(-A_0)$: if $x$ is outside this Minkowski sum, $x+z$ is outside
$C'$ for every $z\in A_0$.  The local finiteness and compact support
make the sum of these contributions locally finite.  Multiplication
by $H$ therefore retains precisely the contribution of $C_0$.
This argument is pointwise and does not discard boundary atoms.
\end{proof}

\begin{lemma}[Rational polygon factors and normalization]
\label{lem:moment-rational-factor}
Suppose the recovered convex body $\widehat C$ has a finite
computable support representation and
$d_H(\widehat C,C_0)\le\xi$.  For every positive rational $\zeta$
one can construct a rational convex polygon $D\supset\widehat C$
with $d_H(D,\widehat C)\le\zeta$.  Put $\tau=\xi+\zeta$ and assume
it is sufficiently small.  Then $|D|\ge v_*/2$ and
\begin{equation}\label{eq:moment-factor-TV}
 \left\|\frac{\one_{C_0}}{|C_0|}\dd x
              -\frac{\one_D}{|D|}\dd x\right\|_{\rm TV}
       \le C\tau.
\end{equation}
All normalized moments
\begin{equation}\label{eq:moment-rational-factor-moments}
 c_\alpha=\frac1{|D|}\int_D(x/R)^\alpha\dd x
\end{equation}
are rational and can be computed exactly by finite arithmetic;
$c_0=1$.  Their discrepancies from the corresponding moments of
the uniform law on $C_0$ are at most $C\tau$, independently of
degree.

Suppose rational estimates $\widehat I_\alpha$ of the occupation
moments have errors at most $\eta$.  Define
\begin{equation}\label{eq:moment-rational-data}
 \widetilde y_0=1,\qquad
 \widetilde y_\alpha=\widehat I_\alpha/|D|
                              \quad(\alpha\ne0).
\end{equation}
If $X$ is uniform on $C_0$ and independent of $Z\sim\mu_0$, then
\begin{equation}\label{eq:moment-normalization-error}
 \left|\widetilde y_\alpha-
                   \E\bigl((X-Z)/R\bigr)^\alpha\right|
       \le C(\eta+\tau).
\end{equation}
\end{lemma}
\begin{proof}
Here is a finite outer-polygon construction.  Take a sufficiently
fine net of rational normal vectors, including the four coordinate
directions.  Compute rational upper bounds for the support of
$\widehat C$ at these vectors, with a prescribed rational error,
and intersect the corresponding half-planes.  The coordinate
directions bound this intersection in a fixed square.  If $q$ is
within $\omega$ of a unit vector $u$, then, for every point $x$ in
that square,
\[
 u\cdot x\le q\cdot x+C\omega,
 \qquad
 |h_{\widehat C}(q)-h_{\widehat C}(u)|\le C\omega.
\]
Consequently, normal mesh and support-value errors at most
$c\zeta$ imply $h_D(u)\le h_{\widehat C}(u)+\zeta$ for every $u$.
The reverse inequality follows from containment.  The vertices of
this bounded intersection are rational.  The finite geometric
representation supplies support values with certified error, so all
operations in this construction terminate.

For convex bodies at Hausdorff distance at most $\tau$, the two
containments in their $\tau$-parallel bodies and the planar
parallel-body formula give
\[
 |C_0\mathbin\triangle D|
       \le (\operatorname{Per}(C_0)+\operatorname{Per}(D))\tau
                    +2\pi\tau^2
       \le 2p_*\tau+2\pi\tau^2.
\]
The relevant perimeters have a uniform bound because the bodies
are convex and lie in a fixed ball.  This also proves
$|D|\ge v_*/2$ for sufficiently small $\tau$.  Directly comparing
the two normalized indicator densities yields
\[
 \left\|\frac{\one_{C_0}}{|C_0|}
              -\frac{\one_D}{|D|}\right\|_1
 \le \frac{2|C_0\mathbin\triangle D|}
                       {\min\{|C_0|,|D|\}},
\]
which proves \eqref{eq:moment-factor-TV}.  A normalized coordinate
monomial is bounded by one on both bodies.  Integration against the
signed difference of their uniform probabilities therefore gives
the degree-independent moment estimate.

Triangulate $D$ using its rational vertices.  An affine map from
the reference triangle converts each monomial integral into a finite
sum of rational multiples of
\[
 \int_{\substack{s,t\ge0\\s+t\le1}}s^p t^q\dd s\dd t
                       =\frac{p!q!}{(p+q+2)!}.
\]
Its Jacobian, coefficients, and the scale $R$ are rational.  This
computes both the area and every integral in
\eqref{eq:moment-rational-factor-moments} exactly.  In particular,
the resulting array is the moment array of the actual uniform
probability on $D$.

Finally, Fubini gives
$I_\alpha=|C_0|\E((X-Z)/R)^\alpha$.  The scale is in the
monomial, so this identity has the physical area $|C_0|$ as its
normalizing factor.  Since $|I_\alpha|\le|C_0|$,
\[
 \left|\frac{\widehat I_\alpha}{|D|}
                    -\frac{I_\alpha}{|C_0|}\right|
 \le \frac{\eta+\bigl||D|-|C_0|\bigr|}{|D|}
 \le C(\eta+\tau).
\]
For degree zero the assigned value is exact.  This proves
\eqref{eq:moment-normalization-error}.
\end{proof}

\subsection{Shared observations and the conditioning bound}
\begin{lemma}[Uniform quadrature for all the occupation moments]
\label{lem:moment-shared-quadrature}
Conditional on the cutoff in Lemma~\ref{lem:moment-protected-cutoff},
let $n\ge1$, $0<\eta<1/4$, and $0<\delta<1/4$.  Fresh observations
with the two fixed commands give rational estimates of all
$I_\alpha$, $|\alpha|\le n$, with simultaneous error at most
$\eta$ and conditional probability at least $1-\delta$, using at
most
\begin{equation}\label{eq:moment-shared-quadrature-cost}
 C(n+1)^2\eta^{-4}
                   \log\frac{C(n+1)}{\eta\delta}
\end{equation}
attempted bits.  The nominal mesh is dyadic and comparable to
$\eta/(n+1)$.  The same empirical occupation values serve every
multi-index.  This conclusion is uniform over all compact launch
probabilities with the stated support bound.
\end{lemma}
\begin{proof}
Tile a fixed square containing the support of $H$, with a fixed
collar, by squares of side $\ell$.  Use one fixed grid corner
$x_Q$ of each square $Q$.  There are at most $C\ell^{-2}$ such
points.  For fixed $z$, put
\[
 w_\alpha(x)=H(x)(x/R)^\alpha,
 \qquad
 f_{\alpha,z}(x)=w_\alpha(x)\one_{C_0}(x+z).
\]
On the square under consideration, the weight is bounded by one
and has Lipschitz constant at most $C(n+1)$.  The union of the
grid squares that meet $\partial(C_0-z)$ has area at most
$C\ell$, uniformly in $z$.  Indeed, it lies in the
$\sqrt2\ell$-neighborhood of this convex boundary, whose area is
bounded by a constant times its perimeter times $\ell$, plus
$C\ell^2$.  On each remaining square the indicator is constant,
and the weight varies by at most $C(n+1)\ell$.  It follows that
\begin{equation}\label{eq:moment-fixed-z-quadrature}
 \left|\ell^2\sum_Q f_{\alpha,z}(x_Q)
                   -\int f_{\alpha,z}(x)\dd x\right|
                       \le C(n+1)\ell.
\end{equation}
Values on a closed obstacle boundary occur only in the boundary
squares already counted.  Thus the estimate holds for every $z$,
and integration against any $\mu_0$ preserves it.  This establishes
the required quadrature bound without regularity of the occupation
field.

At every $x_Q$, use the $2M$ means in the finite-prefix formula.
Estimate each of these means to accuracy $\sigma/(2M)$ with fresh
Bernoulli observations.  The Lipschitz bound for the prefix maximum
then gives $|\widehat v(x_Q)-v(x_Q)|\le\sigma$; truncation to
$[0,1]$ preserves this bound.  Hoeffding's inequality and a union
bound over the at most $2MC\ell^{-2}$ means require at most
\[
 C\ell^{-2}\sigma^{-2}\log\frac{C}{\ell\delta}
\]
attempted bits.  Constants absorb the fixed powers of $M$.
On the resulting event, replacing $v(x_Q)$ by $\widehat v(x_Q)$
in any weighted sum changes it by at most $C\sigma$.
Choose $\ell$ dyadic and comparable to $c\eta/(n+1)$ and set
$\sigma=c\eta$, with $c$ sufficiently small.  Together with
\eqref{eq:moment-fixed-z-quadrature}, this proves the simultaneous
accuracy and \eqref{eq:moment-shared-quadrature-cost}.  All sums
are rational.  The confidence bound is conditional on earlier data
because this stage uses fresh preparations.  The event controlling
the occupation values is common to every monomial, so the number
of moments does not multiply the observation bound.
\end{proof}

\begin{lemma}[An explicit amplification factor]
\label{lem:moment-explicit-conditioning}
Under the hypotheses of Lemma~\ref{lem:two-field-moment-separation},
with moment discrepancy bounded by $b$ through degree $4m$, one has
\begin{equation}\label{eq:moment-explicit-amplification}
 W_1(\mathcal L(Z),\mathcal L(\widetilde Z))
       \le \frac{C}{m}+\Lambda_m b,
 \qquad
 \Lambda_m=C_*m^2 24^{4m}(4m)!,
\end{equation}
where $C_*$ is numerical.  In particular
$\Lambda_m\le(Cm)^{4m+2}$ after increasing the numerical $C$.
\end{lemma}
\begin{proof}
The triangular identity
\eqref{eq:two-field-triangular-moments} gives, for the largest
degree-$k$ discrepancy $D_k$ in the unknown factor,
\[
 D_0=0,\qquad
 D_k\le b(1+2^k)+\sum_{j=0}^{k-1}\binom{k}{j}D_j.
\]
The coefficient of the unknown degree-$k$ moment has modulus one;
all moments of the two individual factors are bounded by one.
Induction gives $D_k\le2b8^k k!$.  To check the induction
quantitatively, the lower-degree sum divided by $2b8^k k!$ is
at most
\[
 \sum_{j=1}^k\frac{8^{-j}}{j!}
            \le\sum_{j=1}^\infty8^{-j}=1/7,
\]
whereas the forcing term divided by the same bound is at most
$1/4$ for $k\ge1$.  The case $b=0$ follows directly from the
triangular identity.

The tensor polynomial construction in the proof of
Lemma~\ref{lem:two-field-moment-separation} approximates any
one-Lipschitz function, after subtraction of its value at zero,
with uniform error at most $C/m$.  Its total degree is at most
$4m$, and the sum of its absolute monomial coefficients is at
most $Cm^2 3^{4m}$.  Applying the preceding moment bound gives
an integration error at most
$Cb m^2 3^{4m}8^{4m}(4m)!$.  Transportation duality proves
\eqref{eq:moment-explicit-amplification}.  Finally
$(4m)!\le(4m)^{4m}$ gives the stated elementary upper bound for
$\Lambda_m$.
\end{proof}

\subsection{Positive reconstruction and its finite output}
\begin{lemma}[A positive rational programme with a feasible comparison law]
\label{lem:moment-positive-programme}
Let $D$ and $c_\alpha$ be as in
Lemma~\ref{lem:moment-rational-factor}, and let $K_A$ be a rational
polygon with $d_H(K_A,A_0)\le\rho$, where $\rho$ is a known
positive rational bound on the geometric accuracy event.  For a
sufficiently small positive dyadic $h$, retain the grid
\begin{equation}\label{eq:moment-padded-grid}
 G=\left\{g\in h\Z^2\cap[-1/2,1/2]^2:
   d_\infty(g,R^{-1}K_A)\le\rho/R+2h\right\}.
\end{equation}
It is nonempty and can be computed by rational arithmetic.  Every
point of $R^{-1}A_0$ has a retained grid point within $\sqrt2h$.

Put $n=4m$, and let the rational numbers $\widetilde y_\alpha$
have errors at most $e_y$ as moments of $Y=(X-Z)/R$.  For a
probability vector $p=(p_g)_{g\in G}$ define
\begin{equation}\label{eq:moment-positive-convolution}
 T_\alpha(p)=\sum_{\beta\le\alpha}\binom\alpha\beta
     c_{\alpha-\beta}(-1)^{|\beta|}
                              \sum_{g\in G}p_g g^\beta.
\end{equation}
The finite programme
\begin{equation}\label{eq:moment-positive-LP}
 \min q\quad\hbox{subject to}\quad
 p_g\ge0,\quad\sum_gp_g=1,\quad q\ge0,\quad
 |T_\alpha(p)-\widetilde y_\alpha|\le q
                                      \quad(|\alpha|\le n)
\end{equation}
has a rational optimizer $p^*$.  If $e_f$ is the variation norm in
\eqref{eq:moment-factor-TV}, its minimum satisfies
\begin{equation}\label{eq:moment-feasible-residual}
 q_*\le e_y+e_f+n\sqrt2h.
\end{equation}
The probability $\widehat\mu_0=\sum_gp_g^*\delta_{Rg}$ obeys
\begin{equation}\label{eq:moment-positive-output-bound}
 W_1(\widehat\mu_0,\mu_0)
 \le R\left\{\frac{C}{m}
            +\Lambda_m(2e_y+e_f+4m\sqrt2h)\right\}.
\end{equation}
Every output atom lies within $C(\rho+Rh)$ of $A_0$ and retains
the canonical coordinates supplied by the geometric stage.
\end{lemma}
\begin{proof}
Round a point $z\in R^{-1}A_0$ to a nearest point of $h\Z^2$,
using a fixed rule in case of a tie.  The rounding distance in
maximum norm is at most $h/2$.  The prior coordinate reserve puts
the rounded point in $[-1/2,1/2]^2$ for small $h$, and
\[
 d_\infty(g,R^{-1}K_A)\le\rho/R+h/2.
\]
It is therefore retained.  Conversely, every retained point has
Euclidean distance at most $C(\rho/R+h)$ from $R^{-1}A_0$.
The rational distance computation in
Lemma~\ref{lem:moment-protected-cutoff} proves that $G$ is a
finite, explicitly decidable set.

Let $U=X/R$, let $\widetilde U$ be the scaled uniform point of
$D$, and let $Z_0=Z/R$.  The exact polygon moments imply
\[
 T_\alpha(p)=\E(\widetilde U-Z_p)^\alpha,
 \qquad \mathcal L(Z_p)=\sum_gp_g\delta_g,
\]
with independent factors.  Thus \eqref{eq:moment-positive-convolution}
has a positive probabilistic realization for every feasible $p$.
For fixed $z$ in the allowed box, the monomial
$(u-z)^\alpha$ is bounded by one.  Replacing $U$ by
$\widetilde U$ changes its expectation, even after averaging in
$z$, by at most $e_f$.  On the same box its gradient with respect
to $z$ has Euclidean norm at most $|\alpha|$.  Quantizing $Z_0$
by the rounding just described therefore changes the convolution
moment by at most $n\sqrt2h$.

The resulting grid law, whose weights need not be rational, is a
comparison point of the real probability simplex.  Its discrepancy
from the rational data is at most $e_y+e_f+n\sqrt2h$, which
proves \eqref{eq:moment-feasible-residual}.  To obtain a rational
optimizer, add the harmless bound
$q\le1+\max_{|\alpha|\le n}|\widetilde y_\alpha|$.
Every $T_\alpha$ has absolute value at most one, so this makes a
nonempty compact rational polytope without changing the minimum.
A linear functional attains its minimum at a vertex, and a vertex
is obtained by solving a nonsingular rational linear system.
The optimum and an optimizing vector are therefore rational and
can be found by finite exact linear programming.  The unknown
comparison law was needed only to bound the optimum.

For the optimizer, the true and fitted convolution moments differ
by at most $q_*+e_y\le2e_y+e_f+n\sqrt2h$.  The two individual
geometric factors have moment discrepancy at most $e_f$.
Lemma~\ref{lem:moment-explicit-conditioning} consequently applies
with $b=2e_y+e_f+n\sqrt2h$.  Scaling transportation distance by
$R$ proves \eqref{eq:moment-positive-output-bound}.  The support
claim follows from the padding estimate already established.
\end{proof}

\begin{corollary}[A sparse positive output]
\label{cor:moment-positive-compression}
The optimizer in Lemma~\ref{lem:moment-positive-programme} can be
chosen with at most
\begin{equation}\label{eq:moment-sparse-output-size}
             \binom{4m+2}{2}=8m^2+6m+1
\end{equation}
nonzero rational weights.  This reduction preserves every fitted
convolution moment through degree $4m$, the objective value, and
the transportation bound.  It uses no further observations.
\end{corollary}
\begin{proof}
There are $d=\binom{4m+2}{2}-1$ nonconstant monomials of total
degree at most $4m$.  Associate to each retained grid point the
rational vector of these $d$ convolution moments together with a
first coordinate equal to one.  If more than $d+1$ weights of a
rational optimizer are positive, the corresponding columns have
a nonzero rational linear dependence $v$.  Its first coordinate
gives $\sum_gv_g=0$, so $v$ has both positive and negative
entries.  Set
\[
       \lambda=\min_{v_g>0}p_g/v_g,
                  \qquad p'_g=p_g-\lambda v_g.
\]
The new weights are nonnegative and rational; their sum and all
the convolution moments are unchanged.  At least one positive
weight vanishes.  Repeating this finite reduction proves the
stated bound.  The hypotheses used in
\eqref{eq:moment-positive-output-bound} are preserved exactly.
\end{proof}

\begin{proposition}[The complete finite factorization budget]
\label{prop:moment-complete-budget}
The construction of Theorem~\ref{thm:two-field-finite-law} may use
the rational factor, shared acquisition, padded programme, and
sparse output above.  Its stated observation, coordinate, and
transportation bounds are unchanged, and its probability output
can have at most $\binom{4m+2}{2}$ atoms.
Conditional on geometric estimates with the stated bounds, the
entire moment construction remains valid for an arbitrary compact
launch probability.
\end{proposition}
\begin{proof}
Let $a_m$ be the dyadic tolerance in
\eqref{eq:two-field-moment-tolerance}.  Reconstruct geometry at
tolerance $c a_m$ with failure allowance $\delta/2$.  Refine the
rational polygons deterministically so that $\xi+\zeta\le Cc a_m$
and $\rho\le Cc a_m$.  Choose the occupation moment tolerance
$\eta=c a_m$ and the normalized mesh $h$ dyadic and comparable
to $c a_m/m$.  The two factor errors in
\eqref{eq:moment-positive-output-bound} satisfy
\[
 e_y\le C a_m,\qquad e_f\le C a_m,
          \qquad 4m\sqrt2h\le C a_m.
\]
Since $\Lambda_m\le(Cm)^{4m+2}$, taking the fixed constant $C_1$
in \eqref{eq:two-field-moment-tolerance} sufficiently large gives
$C\Lambda_m a_m\le1/m$ for every $m\ge2$.  The positive output
therefore has $W_1$ error at most $C/m$.  With
$m=\lceil C_0/\varepsilon\rceil$, this is the asserted
$C\varepsilon$ error.  The geometry was obtained at the finer
tolerance $Ca_m$.

Apply Lemma~\ref{lem:moment-shared-quadrature} with $n=4m$ and
the remaining failure allowance $\delta/2$.  Its conditional
confidence bound is uniform over every outcome in the geometric
accuracy event.  Thus both stages succeed with probability at
least $1-\delta$, and their observation total is
\[
 N_{\rm geom}(c a_m,\delta/2)
          +Cm^2a_m^{-4}\log\frac{Cm}{a_m\delta}.
\]
Substituting the polynomial geometric bound and
$\log(1/a_m)\le C m\log(Cm)$ gives
\eqref{eq:two-field-joint-exponential-cost}.  All occupation
queries use the original vectors $te_1,-te_1$ and their fixed
prefix translates.  Their dyadic coordinate denominators have
logarithms at most $C m\log(Cm)$; the geometric-stage coordinates
have the same bound.  Charging every attempted bit and command
occurrence therefore gives the stated command-description bound.

The candidate grid has at most $C(m/a_m)^2$ points, whereas
Corollary~\ref{cor:moment-positive-compression} bounds the number
of nonzero output atoms by \eqref{eq:moment-sparse-output-size}.
Their locations and weights are rational and are obtained by
finite exact arithmetic.  Construction of the polygons, moment
integration, linear programming, and weight reduction use no
additional observations.  Their arithmetic work and the
representation lengths of the output weights are separate from
the attempted-bit and command-description bounds.  The moment
proofs used compact support, the component gap, and the geometric
factor estimates; they used no density assumption.  For the joint
uniform experiment, the quantitative density priors enter through
$N_{\rm geom}$ exactly as stated in the geometric theorem.
\end{proof}

\subsection*{Source history and assisted analysis}
The technical supplement contains the complete earlier localized,
stationary, calibration, period and directional-germ proofs. The
finite endpoint comparison originates in the earlier internal A2
derivations; the signed-support cancellation and the two-command
finite inverse are proved in this article. The repository records
their derivation history and the point-by-point response to the review.
AI-assisted analysis contributed to these arguments and their
independent checks. The named author is responsible for the proofs
and references. Finite diagnostics and successful compilation provide
reproducibility evidence; they are not certificates of continuum
proofs or of a physical apparatus.
\begin{thebibliography}{99}
\raggedright
\bibitem{AFP}
L. Ambrosio, N. Fusco, and D. Pallara, \emph{Functions of Bounded
Variation and Free Discontinuity Problems}, Oxford Mathematical
Monographs, Clarendon Press, Oxford, 2000.
\bibitem{Altman}
E. Altman, \emph{Constrained Markov Decision Processes},
Chapman \& Hall/CRC, Boca Raton, 1999.
\bibitem{AverkovBianchi}
G. Averkov and G. Bianchi,
\emph{Confirmation of Matheron's conjecture on the covariogram of a planar convex body},
J. Eur. Math. Soc. \textbf{11} (2009), 1187--1202.
\href{https://doi.org/10.4171/JEMS/179}{doi:10.4171/JEMS/179}.
\bibitem{BianchiCross}
G. Bianchi, \emph{The cross covariogram of a pair of polygons determines both
polygons, with a few exceptions}, Adv. Appl. Math. \textbf{42} (2009), 519--544.
\href{https://doi.org/10.1016/j.aam.2008.10.002}{doi:10.1016/j.aam.2008.10.002}.
\bibitem{BianchiLaplace}
G. Bianchi, \emph{The covariogram and Fourier--Laplace transform in $\mathbb C^n$},
Proc. Lond. Math. Soc. (3) \textbf{113} (2016), 1--23.
\href{https://doi.org/10.1112/plms/pdw020}{doi:10.1112/plms/pdw020}.
\bibitem{BGK}
G. Bianchi, R. J. Gardner, and M. Kiderlen,
\emph{Phase retrieval for characteristic functions of convex bodies and
reconstruction from covariograms}, J. Amer. Math. Soc. \textbf{24} (2011), 293--343.
\href{https://doi.org/10.1090/S0894-0347-2010-00683-2}
{doi:10.1090/S0894-0347-2010-00683-2}.
\bibitem{BrunelCaps}
V.-E. Brunel, \emph{Uniform behaviors of random polytopes under the
Hausdorff metric}, Bernoulli \textbf{25} (2019), 1770--1793.
\href{https://doi.org/10.3150/18-BEJ1035}{doi:10.3150/18-BEJ1035}.
\bibitem{BKY}
V.-E. Brunel, J. M. Klusowski, and D. Yang,
\emph{Estimation of convex supports from noisy measurements},
Bernoulli \textbf{27} (2021), 772--793.
\href{https://doi.org/10.3150/20-BEJ1229}{doi:10.3150/20-BEJ1229}.
\bibitem{CGLSupport}
J. Capitao-Miniconi, \'E. Gassiat, and L. Leh\'ericy,
\emph{Support and distribution inference from noisy data}, Bernoulli, forthcoming.
\href{https://arxiv.org/abs/2304.09452}{arXiv:2304.09452}.
Theorem numbering follows the
\href{https://www.e-publications.org/ims/submission/BEJ/user/submissionFile/68730?confirm=a078ae36}
{accepted manuscript}, accessed October 4, 2026.
\bibitem{CN}
R. M. Castro and R. D. Nowak, \emph{Minimax bounds for active learning},
IEEE Trans. Inform. Theory \textbf{54} (2008), 2339--2353.
\href{https://doi.org/10.1109/TIT.2008.920189}{doi:10.1109/TIT.2008.920189}.
\bibitem{DelaigleMeister}
A. Delaigle and A. Meister,
\emph{Nonparametric function estimation under Fourier-oscillating noise},
Statist. Sinica \textbf{21} (2011), 1065--1092.
\href{https://doi.org/10.5705/ss.2009.082}{doi:10.5705/ss.2009.082}.
\bibitem{Galerne}
B. Galerne, \emph{Computation of the perimeter of measurable sets via
their covariogram. Applications to random sets},
Image Anal. Stereol. \textbf{30} (2011), 39--51.
\href{https://doi.org/10.5566/ias.v30.p39-51}
{doi:10.5566/ias.v30.p39-51}.
\bibitem{GLL}
\'E. Gassiat, S. Le Corff, and L. Leh\'ericy,
\emph{Deconvolution with unknown noise distribution is possible for multivariate
signals}, Ann. Statist. \textbf{50} (2022), 303--323.
\href{https://doi.org/10.1214/21-AOS2106}{doi:10.1214/21-AOS2106}.
\bibitem{GK}
P. W. Goldberg and S. Kwek,
\emph{The precision of query points as a resource for learning convex
polytopes with membership queries},
Proceedings of the Thirteenth Annual Conference on Computational
Learning Theory, Morgan Kaufmann, 2000, 225--235.
\bibitem{KiderlenRataj}
M. Kiderlen and J. Rataj, \emph{Dilation volumes of sets of finite perimeter},
Adv. Appl. Probab. \textbf{50} (2018), 1095--1118.
\href{https://doi.org/10.1017/apr.2018.52}{doi:10.1017/apr.2018.52}.
\bibitem{LawlerLimic}
G. F. Lawler and V. Limic, \emph{Random Walk: A Modern Introduction},
Cambridge Studies in Advanced Mathematics \textbf{123}, Cambridge
University Press, 2010.
\bibitem{LCK}
A. Locatelli, A. Carpentier, and S. Kpotufe,
\emph{An adaptive strategy for active learning with smooth decision boundary},
Proceedings of Algorithmic Learning Theory, Proceedings of Machine
Learning Research \textbf{83} (2018), 547--571.
\bibitem{LRSS}
A. K. Louis, M. Riplinger, M. Spiess, and E. Spodarev,
\emph{Inversion algorithms for the spherical Radon and cosine transform},
Inverse Problems \textbf{27} (2011), 035015.
\href{https://doi.org/10.1088/0266-5611/27/3/035015}
{doi:10.1088/0266-5611/27/3/035015}.
\bibitem{MartinezMaure}
Y. Martinez-Maure, \emph{Geometric study of Minkowski differences of
plane convex bodies}, Canad. J. Math. \textbf{58} (2006), 600--624.
\href{https://doi.org/10.4153/CJM-2006-025-x}{doi:10.4153/CJM-2006-025-x}.
\bibitem{MeisterCompact}
A. Meister, \emph{Deconvolving compactly supported densities},
Math. Methods Statist. \textbf{16} (2007), 63--76.
\href{https://doi.org/10.3103/S106653070701005X}{doi:10.3103/S106653070701005X}.
\bibitem{Retained}
Q. Qi, \emph{Scalar collision laws and recognition of periodic dispersing
billiards}, A2 v31, October 4, 2026. Complete reviewed manuscript retained in the repository source archive;
not an independently published article or a required supplementary volume.
\bibitem{RigolletWeed}
P. Rigollet and J. Weed, \emph{Uncoupled isotonic regression via minimum
Wasserstein deconvolution}, Inform. Inference \textbf{8} (2019), 691--717.
\href{https://doi.org/10.1093/imaiai/iaz006}{doi:10.1093/imaiai/iaz006}.
\bibitem{Schneider}
R. Schneider, \emph{Convex Bodies: The Brunn--Minkowski Theory},
second expanded edition, Encyclopedia of Mathematics and its Applications
\textbf{151}, Cambridge University Press, 2014.
\bibitem{Shannon}
C. E. Shannon, \emph{A mathematical theory of communication},
Bell System Tech. J. \textbf{27} (1948), 379--423, 623--656.
\bibitem{Villarrubia}
J. S. Villarrubia, \emph{Algorithms for scanned probe microscope image
simulation, surface reconstruction, and tip estimation},
J. Res. Natl. Inst. Stand. Technol. \textbf{102} (1997), 425--454.
\href{https://doi.org/10.6028/jres.102.030}{doi:10.6028/jres.102.030}.
\bibitem{YCJ}
S. Yan, K. Chaudhuri, and T. Javidi,
\emph{Active learning from imperfect labelers},
Advances in Neural Information Processing Systems \textbf{29} (2016).
\end{thebibliography}

\end{document}
